\documentclass[10pt,letterpaper]{article}

\usepackage[T1]{fontenc}
\usepackage[utf8]{inputenc}
\usepackage{amsmath,amssymb,mathtools,amsthm}
\usepackage{newtxtext,newtxmath}
\usepackage[letterpaper,margin=0.82in,headheight=14pt]{geometry}
\usepackage{microtype}
\usepackage{booktabs,longtable,tabularx,array,multirow}
\usepackage[table]{xcolor}
\usepackage{enumitem}
\usepackage{ragged2e}
\usepackage{needspace}
\usepackage{xurl}
\usepackage{hyperref}
\usepackage{bookmark}
\usepackage[capitalise,nameinlink,noabbrev]{cleveref}
\usepackage{fancyhdr}
\usepackage{lastpage}
\usepackage[most]{tcolorbox}
\usepackage{titlesec}
\usepackage{etoolbox}

\definecolor{navy}{HTML}{17365D}
\definecolor{bluegray}{HTML}{EAF0F6}
\definecolor{green}{HTML}{2F6B4F}
\definecolor{greenbg}{HTML}{EAF4EE}
\definecolor{amber}{HTML}{8A5A00}
\definecolor{amberbg}{HTML}{FFF5DA}
\definecolor{red}{HTML}{8B2E2E}
\definecolor{redbg}{HTML}{FBECEC}
\definecolor{graytext}{HTML}{4F5963}
\definecolor{graybg}{HTML}{F3F5F7}
\definecolor{rulegray}{HTML}{C7CDD4}

\hypersetup{
  colorlinks=true,
  linkcolor=navy,
  citecolor=green,
  urlcolor=navy,
  pdftitle={Literature Synthesis for the Alaoglu--Erdos Integer-Exponent Conjecture},
  pdfauthor={Consolidated from four independent literature surveys},
  pdfsubject={Mathematical literature relevant to integral powers of 2 and 3}
}
\urlstyle{same}

\setlength{\parindent}{0pt}
\setlength{\parskip}{5.2pt plus 1.2pt minus 0.8pt}
\setlength{\emergencystretch}{3em}
\renewcommand{\arraystretch}{1.14}
\setlist[itemize]{leftmargin=1.45em,itemsep=2pt,topsep=3pt}
\setlist[enumerate]{leftmargin=1.65em,itemsep=3pt,topsep=3pt}

\titleformat{\section}{\Large\bfseries\color{navy}}{\thesection}{0.65em}{}
\titleformat{\subsection}{\large\bfseries\color{navy}}{\thesubsection}{0.65em}{}
\titleformat{\subsubsection}{\normalsize\bfseries\color{green}}{\thesubsubsection}{0.6em}{}
\titlespacing*{\section}{0pt}{2.1ex plus 0.8ex minus 0.3ex}{0.8ex}
\titlespacing*{\subsection}{0pt}{1.7ex plus 0.6ex minus 0.2ex}{0.5ex}
\titlespacing*{\subsubsection}{0pt}{1.3ex plus 0.4ex minus 0.2ex}{0.3ex}

\pagestyle{fancy}
\fancyhf{}
\fancyhead[L]{\small\color{graytext}Alaoglu--Erd\H{o}s literature synthesis}
\fancyhead[R]{\small\color{graytext}\nouppercase{\leftmark}}
\fancyfoot[C]{\small\color{graytext}\thepage\ / \pageref*{LastPage}}
\renewcommand{\headrulewidth}{0.35pt}
\renewcommand{\footrulewidth}{0pt}

\newtheorem{proposition}{Proposition}[section]
\newtheorem{question}{Research question}[section]
\newtheorem{criterion}{Success criterion}[section]

\newcommand{\Q}{\mathbb Q}
\newcommand{\R}{\mathbb R}
\newcommand{\Z}{\mathbb Z}
\newcommand{\N}{\mathbb N}
\newcommand{\C}{\mathbb C}
\newcommand{\Gm}{\mathbb G_m}
\newcommand{\Sol}{\mathrm{Sol}}
\newcommand{\Cand}{\mathrm{Cand}}
\newcommand{\ord}{\operatorname{ord}}
\newcommand{\Span}{\operatorname{span}}
\newcommand{\rank}{\operatorname{rank}}
\newcommand{\vpadic}[2]{v_{#1}\!\left(#2\right)}
\newcommand{\BigO}{\mathrm O}
\newcommand{\smallo}{\mathrm o}
\newcommand{\eps}{\varepsilon}
\newcommand{\Rone}{\textsf{R1}}
\newcommand{\Rtwo}{\textsf{R2}}
\newcommand{\Rthree}{\textsf{R3}}
\newcommand{\Rfour}{\textsf{R4}}
\newcommand{\reports}[1]{\textcolor{graytext}{\footnotesize[provenance: #1]}}

\newcommand{\badgeDirect}{\textcolor{green}{\textbf{D}}}
\newcommand{\badgeAdapt}{\textcolor{navy}{\textbf{A}}}
\newcommand{\badgeTemplate}{\textcolor{amber}{\textbf{T}}}
\newcommand{\badgeDiagnostic}{\textcolor{red}{\textbf{X}}}
\newcommand{\badgeCompute}{\textcolor{graytext}{\textbf{C}}}

\newtcolorbox{executivebox}{
  colback=bluegray,colframe=navy,boxrule=0.7pt,arc=1.5mm,
  left=3mm,right=3mm,top=2mm,bottom=2mm,
  before skip=5pt,after skip=7pt
}
\newtcolorbox{takeawaybox}[1][]{
  colback=greenbg,colframe=green,boxrule=0.55pt,arc=1.2mm,
  left=2.7mm,right=2.7mm,top=1.8mm,bottom=1.8mm,
  title={#1},fonttitle=\bfseries,before skip=5pt,after skip=6pt
}
\newtcolorbox{warningbox}[1][]{
  colback=amberbg,colframe=amber,boxrule=0.55pt,arc=1.2mm,
  left=2.7mm,right=2.7mm,top=1.8mm,bottom=1.8mm,
  title={#1},fonttitle=\bfseries,before skip=5pt,after skip=6pt
}
\newtcolorbox{stopbox}[1][]{
  colback=redbg,colframe=red,boxrule=0.55pt,arc=1.2mm,
  left=2.7mm,right=2.7mm,top=1.8mm,bottom=1.8mm,
  title={#1},fonttitle=\bfseries,before skip=5pt,after skip=6pt
}

\newcolumntype{L}[1]{>{\RaggedRight\arraybackslash}p{#1}}
\newcolumntype{C}[1]{>{\Centering\arraybackslash}p{#1}}
\newcolumntype{Y}{>{\RaggedRight\arraybackslash}X}

\begin{document}
\hypersetup{pageanchor=false}

\begin{titlepage}
\thispagestyle{empty}
\vspace*{0.12\textheight}
{\color{navy}\rule{\textwidth}{1.2pt}}\\[1.1em]
{\fontsize{25}{31}\selectfont\bfseries\color{navy}
Literature Synthesis for the\\[0.18em]
Alaoglu--Erd\H{o}s Integer-Exponent Conjecture\par}
\vspace{1.1em}
{\Large\color{graytext}
A critical merge of four independent surveys, aligned with the remaining barriers in
\texttt{alaoglu-erdos-unified-report.tex}\par}
\vspace{1.25em}
{\color{navy}\rule{\textwidth}{0.55pt}}
\vfill
\begin{tabular}{@{}ll@{}}
\textbf{Problem:} & $2^x,3^x\in\Z \Longrightarrow x\in\Z$\\[0.35em]
\textbf{Inputs:} & four literature reports in \texttt{reports/}\\[0.35em]
\textbf{Method:} & deduplication, source correction, barrier mapping, and ranked synthesis\\[0.35em]
\textbf{Date:} & 23 August 2026
\end{tabular}
\vspace*{0.08\textheight}
\end{titlepage}

\hypersetup{pageanchor=true}
\pagenumbering{roman}
\tableofcontents
\clearpage
\pagenumbering{arabic}

\section{Executive synthesis}

\begin{executivebox}
\textbf{Bottom line.}
None of the four surveys contains an off-the-shelf theorem that proves the Alaoglu--Erd\H{o}s
conjecture, and none bypasses the same terminal obstruction isolated in the baseline report:
a nontrivial integer-coefficient relation among products of logarithms of primes.  The merge
nevertheless changes the research portfolio in a useful way.  It identifies two bodies of
literature that are unusually close to the exact arithmetic geometry of a hypothetical
counterexample, and three others that provide concrete mechanisms for creating the extra
``source of size'' missing from the present determinant arguments.
\end{executivebox}

The five strongest clusters are as follows.

\begin{enumerate}
\item \textbf{Integer-valued entire functions on the four-exponentials lattice.}
Pila's two papers on entire functions sharing arguments of integrality, together with
Pila--Rodr\'iguez Villegas on concordant sequences, are the most direct literature found in
any survey.  They place exactly the four algebraic corner values
\(e^{a\alpha},e^{a\beta},e^{b\alpha},e^{b\beta}\) on the semigroup
\(\{i\alpha+j\beta:i,j\geq 0\}\), and give an explicit quadratic-growth threshold below
which a third integer-valued entire function must be algebraically dependent on
\(e^{az}\) and \(e^{bz}\).  For the present problem the threshold is explicit.  The immediate
question is therefore numerical rather than philosophical: can the canonical interpolants
already constructed in the baseline report be placed below that threshold, and if so can
the resulting algebraic dependence be made contradictory?
\reports{\Rfour, \Rthree}

\item \textbf{Arithmetic holonomy and irrationality of products of logarithms.}
Calegari--Dimitrov--Tang prove genuine irrationality and linear-independence results for
certain products \(\log(1+1/m)\log(1+1/n)\), using arithmetic holonomy and an integrated
growth argument.  Their theorem does not include \(\log 2\log 3\), and it must not be quoted
as if it did.  Its importance is methodological: it is a recent successful instance in which
a product-of-logarithms problem acquires a second arithmetic size constraint, precisely the
sort of phenomenon the baseline report lacks.
\reports{\Rthree, \Rfour}

\item \textbf{Multivariate, composite-ideal, and adelic generalized factorials.}
The baseline report convincingly closes the plain one-dimensional, prime-by-prime Bhargava
factorial route.  It does not close the extensions found by the first survey: multivariate
fixed divisors on the actual two-dimensional orbit grid; orderings for composite integers
and arbitrary ideals; adelic regular bases; and matroid/greedoid structure of optimal
ordering sets.  These are the most testable new arithmetic leads because they can be
implemented on finite orbit grids and assessed at the exact asymptotic scale required by the
determinant deficit.
\reports{\Rone}

\item \textbf{Global and \(p\)-adic determinant architecture.}
Heath-Brown, Salberger, Walsh, and recent global-field refinements explain how congruence
coincidence can be converted into an auxiliary algebraic relation rather than merely summed
term by term.  Their theorems apply to algebraic varieties, not to the transcendental arc
\(y=x^{\log_2 3}\), so they are templates rather than plug-ins.  They nevertheless give the
right design discipline for Priority~A: local divisibility should be used to force a global
rank drop before the archimedean estimate is taken.
\reports{\Rthree, \Rfour}

\item \textbf{Conjugate control, house bounds, and adelic capacity.}
Dimitrov's Schinzel--Zassenhaus theorem, Pritsker's variants, and the transfinite-diameter
machinery of DeMarco--Rumely offer a possible route from the Newton--Kummer algebraic integer
to a contradiction.  They become relevant only after one controls enough conjugates or
reciprocity; without that extra step they are comparison theorems, not solutions.  The same
capacity formalism may also prove a no-go theorem showing that a proposed determinant family
is exactly capacity-critical.
\reports{\Rone, \Rfour}
\end{enumerate}

\begin{takeawaybox}[Recommended first moves]
The merged evidence supports four immediate calculations: (i) compute the precise
order-two growth constant of the canonical entire interpolant on \(\Sol\); (ii) compute
multivariate and composite-ideal factorial defects on finite orbit grids and determine their
leading scale; (iii) prototype a determinant whose \(2\)- and \(3\)-adic divisibility is
extracted globally rather than termwise; and (iv) audit the Calegari--Dimitrov--Tang proof to
identify exactly which part of its holonomy construction excludes the pair \((\log 2,\log 3)\).
Each calculation has a clean stop condition and can therefore eliminate a direction without
another large speculative program.
\end{takeawaybox}

\subsection{What the merge rules out}

Several attractive-sounding routes in the supplied reports do not survive contact with the
exact hypotheses.

\begin{itemize}
\item A hypothetical counterexample is necessarily irrational, since rational nonintegral
exponents are ruled out by unique factorization.  Results devoted only to rational \(x\) add
no leverage.
\item The set \(\{(2^t,3^t):t\in\R\}\) is a transcendental analytic curve, not an algebraic
curve in \(\Gm^2\).  Bombieri--Masser--Zannier, Maurin, Zilber--Pink, and arithmetic unlikely
intersection theorems do not apply directly merely because the orbit points are
\(S\)-integral.
\item Classical Baker theory controls linear forms in logarithms with algebraic
coefficients.  It does not control the terminal coefficients
\(m_p\log 3-a_p\log 2\), which are themselves logarithmic.
\item The formal power series \((1-z)^\theta\) with transcendental \(\theta=\log_2 3\) is not
a \(G\)-function in the standard Andr\'e sense, because its coefficients are not algebraic.
Arithmetic algebraization theorems for \(G\)-functions cannot be invoked without first
constructing a genuinely arithmetic replacement.
\item Generic rational-point counting tolerates far more points than the compact orbit
produces.  It must be strengthened to see prime support and synchronized denominator
exponents; effectivity alone is not enough.
\end{itemize}

\subsection{Classification used throughout}

Every lead is marked by one of the following classes.

\begin{center}
\begin{tabularx}{0.96\textwidth}{C{0.08\textwidth}L{0.18\textwidth}Y}
\toprule
Mark & Class & Meaning in this report\\
\midrule
\badgeDirect & Direct & The theorem is formulated in the four-exponentials lattice or in an
essentially identical integer-valued setting.\\
\badgeAdapt & Adaptation & The theorem could apply after one explicit, nontrivial lemma whose
content is stated.\\
\badgeTemplate & Template & The proof architecture is relevant, but the published hypotheses
are materially different.\\
\badgeDiagnostic & Diagnostic & The result identifies a frontier, obstruction, or reason not to
spend effort on a route.\\
\badgeCompute & Computational & The literature supplies algorithms, finite certificates, or
formalization infrastructure rather than an asymptotic theorem.\\
\bottomrule
\end{tabularx}
\end{center}

\clearpage
\section{Baseline, notation, and the exact wall}
\label{sec:baseline}

This synthesis treats the attached unified report as the exclusion baseline.  It therefore
does not re-present the extensive structure theory, determinant calculations, Lean
formalizations, or failed approaches already developed there.  The minimum notation needed
to evaluate the new literature is recorded here.

Set
\[
  \theta=\frac{\log 3}{\log 2}.
\]
A counterexample is a positive irrational \(x\) for which both \(2^x\) and \(3^x\) are
positive integers.  Under the least-generator normalization of the baseline report, let
\(\beta\) be the least positive nonintegral solution and put
\[
  M=2^\beta,\qquad A=3^\beta.
\]
The conditional solution and candidate monoids then have the rigid forms
\[
  \Sol=\N_0+\N_0\beta,
  \qquad
  \Cand=\{2^nM^k:n,k\geq0\}.
\]
The full rectangular orbit grid is
\[
  S_K=\{(2^a3^b,M^aA^b):0\leq a,b<K\}\subset\Z^2,
\]
whereas the compact normalization is
\[
  U_k=\frac{M^k}{2^{\lfloor k\beta\rfloor}}=2^{\{k\beta\}}\in\Z[1/2],
  \qquad
  V_k=\frac{A^k}{3^{\lfloor k\beta\rfloor}}=3^{\{k\beta\}}\in\Z[1/3].
\]
Thus \((U_k,V_k)\) lies on the compact analytic arc
\[
  \mathcal A=\{(2^t,3^t):0\leq t<1\}.
\]

The primitive output equation is
\[
  \log M\,\log 3=\log A\,\log 2.
\]
After prime factorization
\(M=\prod_p p^{m_p}\) and \(A=\prod_p p^{a_p}\), it becomes
\begin{equation}
  D=\sum_p\bigl(m_p\log 3-a_p\log 2\bigr)\log p=0.
  \label{eq:terminal-wall}
\end{equation}
This is an integer-coefficient linear form in products of two prime logarithms.  It is not a
linear form in logarithms with algebraic coefficients, and hence is not covered by Baker's
theory.  The literature is useful only if it does one of the following before reaching
\cref{eq:terminal-wall}:

\begin{enumerate}[label=(\alph*)]
\item creates an additional exact algebraic relation;
\item creates a second, independent lower bound or growth constraint;
\item converts synchronized local divisibility into a global determinant factor;
\item controls a sufficiently rich family of algebraic conjugates; or
\item exploits positivity, integrality, or prime support in a way unavailable for a general
four-exponentials matrix.
\end{enumerate}

\subsection{Barrier model}

The four surveys use different vocabularies.  The following common barrier model is used in
the synthesis and in the crosswalk tables.

\begingroup
\small
\setlength{\LTleft}{0pt}
\setlength{\LTright}{0pt}
\begin{longtable}{@{}L{0.12\textwidth}L{0.82\textwidth}@{}}
\toprule
Barrier & Exact content\\
\midrule
\endfirsthead
\toprule
Barrier & Exact content\\
\midrule
\endhead
\bottomrule
\endfoot
\textbf{B1: synchronized denominators} & A useful estimate must depend jointly on the
\(2\)- and \(3\)-adic data, or on their common floor parameter, rather than separately
optimizing two incompatible local problems.\\
\addlinespace
\textbf{B2: rank collapse} & A counterexample produces a multiplicative semigroup of rank at
most two.  Existing rank theorems do not force rank one in the balanced \(2\times2\) case.\\
\addlinespace
\textbf{B3: log-product wall} & The final identity is \cref{eq:terminal-wall}.  A classical
linear-form estimate loses algebraicity of the coefficients exactly here.\\
\addlinespace
\textbf{B4: \(p\)-adic amplification} & A determinant needs local divisibility on the same
leading scale as its archimedean size.  Factoring rowwise minima or summing termwise tropical
bounds is too weak.\\
\addlinespace
\textbf{B5: conjugate compatibility} & Kummer conjugates of an algebraic radical do not
transport the selected real logarithm relation equivariantly.  Norms alone recover known
identities.\\
\addlinespace
\textbf{B6: effective height gap} & Available irrationality measures are polynomial in
height, while the integral gap created by the candidate equations is exponentially small in
the natural parameter.\\
\addlinespace
\textbf{B7: support-sensitive counting} & The compact arc carries only logarithmically many
ordinary rational points, so generic o-minimal bounds are vacuous.  One must count points
with disjoint prescribed prime supports and synchronized exponents.\\
\addlinespace
\textbf{B8: full-grid sparsity} & Ordinary interpolation on a \(K^2\)-point grid uses one
monomial too many, or loses a comparable factor in degree.  The desired relation must save
support or exploit arithmetic not measured by support.\\
\addlinespace
\textbf{B9: local--global optimization} & Separately optimal local orderings cannot simply be
summed.  The invariant must select rows, columns, and places jointly.\\
\end{longtable}
\endgroup

\subsection{Crosswalk to the baseline research priorities}

\begin{center}
\small
\begin{tabularx}{\textwidth}{L{0.075\textwidth}L{0.27\textwidth}Y}
\toprule
Priority & Baseline target & Most relevant new literature\\
\midrule
A & \(p\)-adic amplification of semigroup determinants & Heath-Brown--Salberger--Walsh;
multivariate/composite orderings; matroid defects; recent global determinant methods.\\
B & Full \((n,k)\)-grid and sparse interpolation & Prasad--Rajkumar--Reddy; Evrard;
multivariate fixed divisors; Pila's geometric/arithmetic postulation.\\
C & Denominator-sensitive compact-arc counting & Campagna--Dill--Wilms; effective Pfaffian
counting as a benchmark; analytic gcd bounds; composite ideal orderings.\\
D & Mixed archimedean--radical invariants & Dimitrov; Pritsker; DeMarco--Rumely; Bost slopes;
adelic capacity.\\
E & Sharpen determinant hypotheses & Bombieri--Pila; global determinant method; Pad\'e and
potential theory; Habsieger's alternants.\\
F & Force a third algebraic exponential & Pila's third entire function; complex-power
geometry only as a speculative existence framework.\\
G & Five/sharp six exponentials & Dasgupta's rank survey; Roy's parametric geometry;
Waldschmidt--Masser frontier.\\
H & Larger effective/verified regions & Simons--de Weger; Pad\'e/Hankel methods; DSLean-style
certificate tooling.\\
I & Exact computational reconnaissance & Algorithms for generalized orderings; matroid
intersection; finite determinant searches; digit/carry certificates.\\
J & Colossally abundant interface & Caveney--Nicolas--Sondow and the original
Alaoglu--Erd\H{o}s application.\\
\bottomrule
\end{tabularx}
\end{center}

\section{Synthesis protocol and correction audit}
\label{sec:audit}

\subsection{How the four reports were merged}

The reports were not concatenated.  Their claims were normalized into a common record with
five fields: theorem actually proved, hypotheses, proposed point of attachment to the
baseline argument, missing lemma, and expected asymptotic scale.  Duplicate references were
combined; incompatible assessments were resolved in favor of the narrower statement; and
bibliographic metadata was corrected where the supplied reports disagreed with the cited
source.

The provenance tags used below have the following meaning:
\(\Rone\) is the generalized-factorial and adelic-ordering survey;
\(\Rtwo\) is the broad transcendence/rank/dynamics survey;
\(\Rthree\) is the products-of-logs and determinant-method dossier; and
\(\Rfour\) is the integer-valued-entire-function, house-bound, and arithmetic-algebraization
survey.  Provenance records discovery, not endorsement.

\subsection{Substantive corrections}

\begingroup
\footnotesize
\setlength{\LTleft}{0pt}
\setlength{\LTright}{0pt}
\begin{longtable}{@{}L{0.21\textwidth}L{0.34\textwidth}L{0.39\textwidth}@{}}
\toprule
Item & Claim appearing in a supplied report & Corrected assessment used here\\
\midrule
\endfirsthead
\toprule
Item & Claim appearing in a supplied report & Corrected assessment used here\\
\midrule
\endhead
\bottomrule
\endfoot
Rational exponents & Results on rational \(x\) were presented as progress toward the
conjecture. & Redundant.  If \(x=p/q\) in lowest terms and \(2^x\in\Z\), unique factorization
already forces \(q=1\).\\
\addlinespace
Algebraic-curve methods & BMZ, Maurin, or Zilber--Pink were described as directly applicable
to the orbit. & The orbit lies on a transcendental analytic arc.  These theorems require an
algebraic subvariety or an arithmetic family satisfying additional algebraic hypotheses.\\
\addlinespace
Calegari--Dimitrov--Tang & Their product-of-logarithms theorem was suggested as nearly covering
\(\log 2\log 3\). & It proves a nearby but different family involving
\(\log(1+1/m)\log(1+1/n)\), under a strong near-equality condition.  It is a method lead, not
a specialization.\\
\addlinespace
\(G\)-functions & \((1-z)^\theta\), \(\theta=\log_2 3\), was called a \(G\)-function. & False
in the standard arithmetic definition: the Taylor coefficients are generally
transcendental.  A separate arithmetic holonomic object would have to be constructed.\\
\addlinespace
Dimitrov's theorem & The house bound was stated without keeping track of reciprocal
orientation. & Dimitrov's theorem is naturally a statement about roots of an integer
polynomial with \(P(0)=1\); using it for a Newton--Kummer algebraic integer requires a
reciprocal polynomial or controlled inverse/conjugate relation.\\
\addlinespace
Adelic regular bases & Applicability to the adelic closure of \(\langle2,3\rangle\) was stated
as essentially automatic. & It is a concrete verification problem.  The residual-cardinality
and compactness hypotheses must be checked place by place, with \(2\) and \(3\) treated
separately.\\
\addlinespace
Composite \(6\)-ordering & A \(6\)-ordering was described as automatically coupling the two
prime directions. & It couples them only through
\(v_{(6)}(n)=\min(v_2(n),v_3(n))\), which may discard rather than exploit an imbalance.  It is
a testable invariant, not a guaranteed gain.\\
\addlinespace
Integer-valued o-minimal functions & arXiv:2404.10737 was attributed to Bhardwaj--van den
Dries. & The paper is by Neer Bhardwaj, Raymond McCulloch, Nandagopal Ramachandran, and
Katharine Woo.\\
\addlinespace
Pila--Rodr\'iguez Villegas & The paper was assigned to the wrong journal in one report. &
\emph{Concordant sequences and integral-valued entire functions}, \emph{Acta Arith.} 88
(1999), 239--268.\\
\addlinespace
Pila's postulation paper & One report gave the wrong year. & \emph{Geometric and arithmetic
postulation of the exponential function}, \emph{J. Austral. Math. Soc. Ser. A} 54 (1993),
111--127.\\
\addlinespace
Pasten--Wang & The venue was listed as \emph{Trans. Amer. Math. Soc.} & \emph{GCD bounds for
analytic functions}, \emph{Int. Math. Res. Not.} 2017, no. 1, 47--95.\\
\addlinespace
Complex powers & Exponential-algebraic closedness was described as a way to manufacture the
missing independent column. & That is speculative and directionally opposite to the needed
rigidity.  The results prove existence/tameness, not arithmetic independence of a prescribed
integral point.\\
\addlinespace
Symbolic dynamics & Chaos or digit complexity was presented as evidence against a
counterexample. & At most heuristic.  Complexity becomes useful only after deriving an
actual automatic, morphic, or finite-state sequence from the counterexample equations.\\
\end{longtable}
\endgroup

\begin{warningbox}[Interpretive rule]
A paper can be highly relevant without its theorem applying.  Throughout the report,
``template'' means that the proof contains a mechanism worth transplanting, not that the
conclusion may be quoted for the Alaoglu--Erd\H{o}s orbit.  This distinction is especially
important for determinant methods, unlikely intersections, \(G\)-function algebraization,
and complex-power geometry.
\end{warningbox}

\section{Ranked research portfolio}
\label{sec:ranking}

The ranking below measures expected value for the particular barriers in
\cref{sec:baseline}, not prestige or general mathematical importance.  A high rank requires
both structural proximity and a falsifiable next calculation.

\begingroup
\footnotesize
\setlength{\LTleft}{0pt}
\setlength{\LTright}{0pt}
\setlength{\tabcolsep}{3.5pt}
\begin{longtable}{@{}C{0.04\textwidth}L{0.255\textwidth}C{0.065\textwidth}L{0.12\textwidth}L{0.44\textwidth}@{}}
\toprule
Rank & Lead & Class & Barriers & Decisive next step\\
\midrule
\endfirsthead
\toprule
Rank & Lead & Class & Barriers & Decisive next step\\
\midrule
\endhead
\bottomrule
\endfoot
1 & Pila, \emph{Entire functions sharing arguments of integrality I--II};
Pila--Rodr\'iguez Villegas & \badgeDirect & B3, B8 & Compute the order-two type of a
canonical entire interpolant integer-valued on \(\Sol\); compare it with
\(1/(832\,\beta^2\log2\log3)\); then classify the algebraic dependence forced below the
threshold.\\
\addlinespace
2 & Prasad--Rajkumar--Reddy and Evrard: multivariate generalized factorials/fixed divisors &
\badgeAdapt & B1, B4, B8, B9 & Apply the sharp multidegree or total-degree factorial to the
actual set \(S_K\subset\Z^2\).  Determine whether the mixed fixed divisor has a positive
surplus on the leading determinant scale.\\
\addlinespace
3 & Calegari--Dimitrov--Tang: arithmetic holonomy and products of logarithms &
\badgeTemplate & B3, B6 & Reconstruct the exact holonomic family, denominator estimates, and
integrated-growth inequality; identify whether a deformation or logarithmic decomposition can
include \((\log2,\log3)\) without destroying arithmeticity.\\
\addlinespace
4 & Lagarias--Yangjit: composite \(b\)-orderings and arbitrary-ideal orderings &
\badgeAdapt & B1, B4, B9 & Compute orderings for \((2),(3),(6)\) and a small family
\((2^u3^v)\) on candidate grids; compare the joint factorial with the sum of separate local
optima at normalized leading order.\\
\addlinespace
5 & Grinberg--Petrov; Krachun--Mirgalimova: matroid/greedoid structure &
\badgeAdapt & B1, B9 & Define the minimum simultaneous \(2\)-/\(3\)-adic defect as a weighted
matroid-intersection problem; prove or disprove a leading-scale lower bound.\\
\addlinespace
6 & Heath-Brown--Salberger--Walsh and global-field determinant refinements &
\badgeTemplate & B4, B8, B9 & Replace termwise valuation bounds by a finite-point auxiliary
hypersurface construction.  The prototype succeeds only if local congruence forces a global
rank drop before the transcendental arc enters.\\
\addlinespace
7 & Dimitrov; Pritsker: house bounds and capacity for algebraic integers &
\badgeAdapt & B5, B6 & Produce a reciprocal or conjugate-stable Newton--Kummer polynomial and
show that all relevant conjugates lie in the region where the house theorem contradicts the
archimedean construction.\\
\addlinespace
8 & Chabert--Peruginelli: adelic orderings and regular bases &
\badgeAdapt & B1, B9 & Verify the hypotheses for the adelic closure of
\(\langle2,3\rangle\), compute the regular basis degrees, and test whether the global basis
contains a mixed term absent from primewise factorials.\\
\addlinespace
9 & DeMarco--Rumely; Rumely: transfinite diameter and adelic capacity &
\badgeTemplate & B5, B9 & Express the orbit determinant family as a capacity problem.  Either
obtain a strict product-formula inequality or prove that the family is capacity-critical and
cannot close the missing logarithm.\\
\addlinespace
10 & Habsieger: alternants related to four exponentials &
\badgeDiagnostic & B3, B8 & Translate the current semigroup determinant into Habsieger's
alternant language and isolate which positivity/divisibility condition would strengthen the
known obstruction.\\
\addlinespace
11 & Bost slopes; Calegari--Dimitrov--Tang unbounded denominators &
\badgeTemplate & B3, B6, B9 & Construct an arithmetic differential or holonomic object from
the candidate lattice whose denominator growth violates a slope inequality.  Without such an
object, the theory has no input.\\
\addlinespace
12 & Schmidt--Summerer and Roy parametric geometry of numbers; shifted-log Pad\'e systems &
\badgeTemplate & B3, B6 & Encode the two simultaneous approximation problems as a combined
successive-minima trajectory and locate an impossible slope pattern forced by exact
integrality.\\
\addlinespace
13 & Simons--de Weger; explicit Pad\'e/Hankel estimates; DSLean-style certificates &
\badgeCompute & B6 & Turn existing irrationality estimates for \(\log2/\log3\) into compact
kernel-replayable exclusion certificates over substantially larger finite ranges.\\
\addlinespace
14 & Adamczewski--Faverjon; function-field analogues; digit/carry literature &
\badgeDiagnostic & B2, B7 & Derive an actual Mahler, automatic, or finite-state structure from
the counterexample.  If no such structure emerges, these remain analogies.\\
\addlinespace
15 & Caveney--Nicolas--Sondow and the colossally abundant interface &
\badgeCompute & downstream & Formalize the implication from the conjecture to primality of
successive colossally abundant quotients, after the core arithmetic modules stabilize.\\
\end{longtable}
\endgroup

\begin{takeawaybox}[Portfolio interpretation]
Ranks 1--6 are active research programs with a finite first experiment.  Ranks 7--12 require
a new structural lemma before their literature can engage.  Ranks 13--15 are valuable for
verified finite progress, diagnostics, and the historical application, but they do not by
themselves threaten the infinite conjecture.
\end{takeawaybox}

\section{Integer-valued entire functions on the solution lattice}
\label{sec:pila}

\subsection{Why this is the closest published program}

Suppose \(a,b,\alpha,\gamma>0\) and the four corner values
\[
  e^{a\alpha},\quad e^{a\gamma},\quad e^{b\alpha},\quad e^{b\gamma}
\]
are integers.  Then the two exponentials
\[
  f_1(z)=e^{az},\qquad f_2(z)=e^{bz}
\]
are integer-valued on the additive semigroup
\[
  X_{\alpha,\gamma}=\{i\alpha+j\gamma:i,j\in\N_0\}.
\]
This is exactly the setup studied by Pila in
\emph{Entire functions sharing arguments of integrality I--II}
\cite{PilaSharingI,PilaSharingII}.  In the Alaoglu--Erd\H{o}s normalization, take
\[
  a=\log2,\qquad b=\log3,\qquad \alpha=1,\qquad \gamma=\beta.
\]
Then
\[
  e^{a\alpha}=2,
  \quad e^{a\gamma}=M,
  \quad e^{b\alpha}=3,
  \quad e^{b\gamma}=A,
\]
and \(X_{1,\beta}=\Sol\).  No change of category is required: the published theorem is
already about the exact additive semigroup forced by a counterexample.

One of Pila's order-two results implies, in this specialization, that an entire function
\(g\) integer-valued on \(\Sol\) and satisfying a sufficiently small quadratic-growth bound,
for example
\begin{equation}
  \limsup_{r\to\infty}\frac{\log M(g,r)}{r^2}
  \leq \frac{1}{832\,\beta^2\log2\log3},
  \label{eq:pila-threshold}
\end{equation}
under the corresponding hypotheses of the theorem, is algebraically dependent with
\(2^z\) and \(3^z\).  Here \(M(g,r)=\max_{|z|\leq r}|g(z)|\).  The constant is not expected
to be optimal; its value is that it turns an open-ended interpolation idea into a measured
target.

\begin{takeawaybox}[Direct numerical test]
For every canonical interpolant already present in the baseline program---Newton,
\(q\)-Newton, Stieltjes, or a minimal-growth interpolation envelope---compute or rigorously
bound
\[
  \tau_2(g):=\limsup_{r\to\infty}\frac{\log M(g,r)}{r^2}.
\]
The first gate is whether \(\tau_2(g)\) can be smaller than the right side of
\cref{eq:pila-threshold}.  This calculation is meaningful even if it fails: the excess
factor measures the same quantitative deficit that appears in the determinant hierarchy.
\end{takeawaybox}

\subsection{What algebraic dependence would and would not give}

Pila's conclusion is not by itself the Alaoglu--Erd\H{o}s conjecture.  It produces a
nonzero polynomial
\[
  P(X,Y,Z)\in\Z[X,Y,Z]
\]
such that
\[
  P(2^z,3^z,g(z))\equiv0.
\]
To obtain a contradiction, one needs one of the following stronger outcomes.

\begin{enumerate}
\item \textbf{Designed transcendence over the exponential field.}
Construct \(g\) in a way that proves \(g\) is transcendental over
\(\C(2^z,3^z)\).  Then Pila's dependence is impossible.

\item \textbf{Classification of low-growth algebraic branches.}
Show that every entire function algebraic over \(\C(2^z,3^z)\) and having the relevant
growth belongs to a small ring, ideally \(\C[2^z,3^z]\) or a finite algebraic extension with
controlled ramification.  Integer values on \(\Sol\), together with normalization conditions,
may then exclude the resulting forms.

\item \textbf{A third algebraically independent interpolation condition.}
Choose \(g\) to encode a quantity not recoverable from \(2^z\) and \(3^z\), such as a
carefully normalized finite-difference cocycle or a denominator-support statistic.  This is
the closest entire-function analogue of Priority~F.
\end{enumerate}

The second route deserves an explicit differential-algebra audit.  Since \(2^z\) and
\(3^z\) solve first-order linear differential equations, an algebraic relation involving an
entire \(g\) may impose strong ramification and growth restrictions.  Those restrictions are
not supplied by Pila's theorem, but they are potentially simpler than proving a new
four-exponentials statement from scratch.

\subsection{Concordant sequences and the minimal-growth envelope}

Pila--Rodr\'iguez Villegas construct minimal-growth transcendental entire functions
integer-valued on broad classes of diffuse concordant sequences
\cite{PilaRodriguezVillegas}.  Geometric progressions are central examples, and the older
Gel'fond--Fukasawa--Gramain theory gives sharp comparison scales for functions integer-valued
on \(\{a^n\}\) \cite{Gramain1981,WaldschmidtIntegerValued}.  There are two distinct node sets
in the present problem:

\[
  \Sol=\N_0+\N_0\beta
  \qquad\text{and}\qquad
  \Cand=\{2^nM^k:n,k\geq0\}.
\]
The first is additively two-generated but lies on the real line; the second is a
multiplicative rank-two set whose logarithms are exactly \(\Sol\log2\).  The concordant
sequence machinery should be tested on both representations, but its hypotheses must not be
assumed.  In particular, one should verify:

\begin{itemize}
\item whether an increasing enumeration of \(\Cand\) is diffuse and concordant in the
technical sense of the paper;
\item whether the envelope construction preserves integrality on the full semigroup rather
than only on the two coordinate progressions;
\item how the envelope's growth changes when the two generators have irrational logarithmic
ratio; and
\item whether the minimal envelope is visibly outside the algebraic closure of
\(\C(2^z,3^z)\).
\end{itemize}

A promising computational approach is to construct finite Newton products over the first
\(N\) nodes of \(\Sol\) and \(\Cand\), normalize them by their fixed divisors, and estimate
the logarithmic potential of the empirical node measure.  If the limit exists, it should
identify the true growth constant before a full concordance proof is attempted.

\subsection{A concrete work package}

\begin{description}[style=nextline]
\item[Input.]
The canonical interpolants and node-ordering machinery already developed in the baseline
report; Pila's exact hypotheses and constants; the concordant-sequence envelope.

\item[Deliverables.]
A rigorous upper and lower bound for \(\tau_2(g)\); a proof or disproof of concordance for the
chosen enumeration; and a classification lemma for entire functions algebraic over
\(\C(2^z,3^z)\) in the attained growth class.

\item[Success condition.]
Either \(\tau_2(g)\) lies below Pila's threshold and the resulting algebraic dependence is
excluded, or the comparison reveals a new structural constraint on any counterexample.

\item[Stop condition.]
If every natural integer-valued interpolant has a rigorously established type exceeding the
threshold by a fixed factor and no algebraic-independence design is available, this route has
quantified rather than removed the one-logarithm deficit and should be paused.
\end{description}

\section{Arithmetic holonomy and products of logarithms}
\label{sec:cdt}

\subsection{The exact result and its boundary}

Calegari--Dimitrov--Tang prove new linear-independence statements involving products of two
logarithms in their work on \(1,\zeta(2)\), and \(L(2,\chi_{-3})\)
\cite{CDTProducts}.  One consequence treats
\[
  \log\left(1+\frac1m\right)
  \log\left(1+\frac1n\right)
\]
for nonexceptional integers \(m,n\) satisfying a stringent near-equality condition.  In a
representative form, the product is irrational, and when \(m\neq n\) it participates with
the two individual logarithms in a \(\Q\)-linearly independent set.

The pair \((\log2,\log3)\) is not obtained by choosing \(m,n\) in this theorem.  Indeed
\(\log2=\log(1+1)\) corresponds to an excluded boundary value, while \(\log3\) is not a
single logarithm of the required neighboring form.  Nor does a formal decomposition such as
\[
  \log3=\log2+\log\frac32
\]
solve the problem: it introduces cross-products and changes the arithmetic holonomic family.
The theorem must therefore be treated as evidence that product-logarithm irrationality can
sometimes be proved, not as a near-complete special case of the present conjecture.

\subsection{Why the method matters}

The relevant feature of the proof is the creation of two independent controls on the same
auxiliary object:

\begin{enumerate}
\item a denominator or integrality estimate coming from arithmetic holonomy; and
\item an archimedean integrated-growth estimate strong enough to force a nonzero linear form
to be too small.
\end{enumerate}

This is conceptually parallel to the baseline determinant program, where integrality gives
\(|D|\geq1\) but the analytic estimate misses by one logarithm.  The CDT machinery succeeds
in a family because holonomy supplies more structure than arbitrary logarithms possess.  The
research question is therefore not ``can their theorem be substituted?'' but ``can the
Alaoglu--Erd\H{o}s lattice be encoded by an arithmetic holonomic family with comparable
local and global controls?''

\subsection{Four attachment experiments}

\subsubsection{Deformation through neighboring logarithms}

Search for an identity expressing the terminal quantity as a specialization or limit of
periods in the CDT family.  A useful deformation must preserve all three properties:

\[
  \text{arithmetic coefficients},\qquad
  \text{uniform denominator control},\qquad
  \text{uniform nondegenerate holonomy}.
\]
A mere analytic limit is insufficient, because irrationality and linear independence are not
closed under specialization.  The key output of this experiment is a proof that the
necessary estimates remain uniform, or a precise statement of where they blow up as the
parameters approach the desired boundary.

\subsubsection{Telescoping decompositions}

Since
\[
  \log q=\sum_j \log\left(1+\frac{1}{m_j}\right)
\]
can sometimes be achieved by telescoping products of neighboring rationals, one can ask
whether both \(\log2\) and \(\log3\) admit decompositions using a common near-diagonal family.
The obstacle is combinatorial explosion of cross-products.  A successful decomposition
would have to place all cross-terms inside a single holonomic period matrix, not simply
invoke pairwise irrationality statements.

\subsubsection{Holonomic encoding of the semigroup determinant}

The determinant entries
\[
  (2^nM^k)^a(3^nA^k)^b
\]
are hypergeometric in each integer index.  One should test whether selected minors, their
finite differences, or their generating series satisfy a Picard--Fuchs or holonomic system
with arithmetic coefficients independent of the hypothetical transcendental parameter
\(\beta\).  The last qualification is essential: a differential equation whose coefficients
contain \(M\), \(A\), or \(\beta\) transcendently has not gained arithmeticity.

\subsubsection{Unbounded denominators as a contradiction}

The CDT proof of the unbounded denominators conjecture for modular forms
\cite{CDTUnbounded} belongs to a broader slope/algebraization tradition.  If a natural series
attached to the candidate grid has bounded denominators because all its sampled values are
integral, while arithmetic holonomy forces unbounded denominators unless the series is
algebraic, then the candidate equations may force algebraicity.  The missing object must be
specified.  The formal binomial series \((1-z)^\theta\) is not suitable because its
coefficients are not algebraic.

\begin{stopbox}[Stop condition for the holonomy route]
Pause the program if every plausible generating series either has transcendental
coefficients from the outset, loses uniform denominator bounds under specialization, or has
a holonomy group too small to encode the product \(\log2\log3\).  ``Holonomic'' without
arithmetic coefficients is not progress on \cref{eq:terminal-wall}.
\end{stopbox}

\subsection{Bost's slope method and arithmetic algebraization}

Bost's algebraicity criteria for formal-analytic leaves and later arithmetic slope methods
provide the general local--global language behind these ideas \cite{Bost2001,BostCharles}.
An evaluation map on sections has an archimedean operator norm and nonarchimedean denominator
bounds; the slope inequality constrains how often it can vanish without forcing algebraicity.
For the present problem, an implementation would require:

\begin{itemize}
\item an algebraic ambient variety defined over a number field;
\item a formal or analytic leaf whose Taylor coefficients are arithmetic;
\item evaluation points or jets derived from \(\Sol\) or \(S_K\); and
\item a global height inequality in which the synchronized \(2\)- and \(3\)-adic gains are
visible before separate optimization destroys them.
\end{itemize}

The baseline interpolants currently satisfy the analytic part of this template but not the
arithmetic-leaf part.  Establishing an arithmetic differential equation is therefore the
first gate, not a technical afterthought.

\clearpage
\section{Mixed generalized factorials, orderings, and fixed divisors}
\label{sec:factorials}

\subsection{What the baseline report has already excluded}

The baseline report develops ordinary Bhargava \(p\)-orderings for the scalar semigroup
\(\langle2,3\rangle\), computes the separate local factorial behavior, and shows that the
locally optimal orderings at \(2\) and \(3\) are incompatible.  That work rules out the
strategy
\[
  \text{choose one scalar ordering}\;\longrightarrow\;
  \text{sum independent \(p\)-adic factorial gains}\;\longrightarrow\;
  \text{beat the archimedean determinant}.
\]
It does not rule out theories whose basic object is already multivariate, composite-ideal, or
adelic.  This distinction is the principal contribution of \Rone.

\subsection{Multivariate fixed divisors on the actual orbit grid}

Evrard introduced multivariate analogues of Bhargava factorials organized by total degree
\cite{Evrard2012}.  Prasad--Rajkumar--Reddy develop sharper fixed-divisor formulas for
polynomials in several variables with multidegree and degree constraints
\cite{PrasadRajkumarReddy}.  Instead of replacing the counterexample by the one-dimensional
set \(\{2^a3^b\}\), one can feed the theory the actual integral point set
\[
  S_K=\{(2^a3^b,M^aA^b):0\leq a,b<K\}\subset\Z^2.
\]

Let \(\Lambda\subset\N_0^2\) be a finite lower set of exponent pairs, and consider the
evaluation matrix
\[
  E_{(a,b),(i,j)}=(2^a3^b)^i(M^aA^b)^j,
  \qquad (i,j)\in\Lambda.
\]
The relevant generalized factorial controls the fixed divisor of every primitive polynomial
\[
  P(X,Y)=\sum_{(i,j)\in\Lambda}c_{ij}X^iY^j
\]
on \(S_K\), or equivalently the unavoidable divisibility of suitable evaluation
minors.  This retains three pieces of information discarded by scalarization:

\begin{itemize}
\item the two coordinate functions have different prime supports at the compact level;
\item the same pair \((a,b)\) controls both coordinates;
\item the Newton polytope \(\Lambda\) can be optimized against the triangular regions used by
the baseline sparse-interpolation program.
\end{itemize}

The theorem itself does not promise a gain.  The decisive statistic is a \emph{mixed fixed-
divisor surplus}.  One useful operational definition is
\begin{equation}
  G_{p,K}(\Lambda)
  :=\min_{T}\,v_p\!\det E_T-L^{\mathrm{flat}}_{p,K}(\Lambda),
  \label{eq:mixed-gain}
\end{equation}
where \(T\) ranges over full-rank choices of rows from \(S_K\) and
\(L^{\mathrm{flat}}_{p,K}(\Lambda)\) is the best lower bound obtainable after replacing the
row by a scalar magnitude and applying the already-analyzed one-dimensional ordering.  The
normalization can be adjusted to the determinant family, but the comparison must be made at
the same cardinality and Newton polytope.

\begin{criterion}
The multivariate route is useful only if
\[
  G_{2,K}(\Lambda_K)\log2+G_{3,K}(\Lambda_K)\log3
\]
has a positive contribution on the leading scale of the archimedean determinant deficit.
An \(O(K)\) boundary correction cannot repair a \(\Theta(K^2)\) loss; in a determinant with
\(N\asymp K^2\) rows, the correct comparison may instead be \(N^2\), depending on the chosen
monomial family.  The scale must be fixed before numerical evidence is interpreted.
\end{criterion}

\subsubsection{Finite experiment}

For \(K\leq 10\) and a collection of lower sets \(\Lambda\)---rectangles, total-degree
triangles, and the baseline boundary-dilation shapes---compute:

\begin{enumerate}
\item exact Smith normal forms of the evaluation matrices;
\item the \(2\)- and \(3\)-adic valuations of all maximal minors or a certified minimum;
\item the multivariate factorial predicted by the fixed-divisor theorem; and
\item the normalized gap from the scalar primewise benchmark.
\end{enumerate}

The Smith profile is more informative than one determinant: a systematic block of invariant
factors would show that the gain survives row/column choices, whereas one exceptional minor
may be a finite accident.

\subsection{Composite \texorpdfstring{\(b\)}{b}-orderings and arbitrary-ideal orderings}

Lagarias--Yangjit extend Bhargava's construction from prime ideals to composite integers and,
more recently, to arbitrary ideals of Dedekind domains, including \(r\)-removed and order-
\(h\) variants \cite{LagariasYangjit2024,LagariasYangjit2025}.  For an integer \(b>1\), the
basic divisibility order is
\[
  v_b(n)=\max\{e\geq0:b^e\mid n\}.
\]
In particular,
\[
  v_6(n)=\min\{v_2(n),v_3(n)\}.
\]
A \(6\)-ordering therefore detects simultaneous divisibility of a difference by both primes.
This is genuinely different from selecting a \(2\)-ordering and a \(3\)-ordering
independently, but it can also be weaker: a large valuation at one prime is ignored when the
other is small.

The right test family is not just \(b=6\).  Use
\[
  \mathcal B_H=\{2^u3^v:0\leq u,v\leq H,\ (u,v)\neq(0,0)\}
\]
or, in ideal notation, the collection \(\{(2^u3^v)\}\).  Each slope \(u/v\) probes a
different compromise between the two local directions.  A weighted envelope over this
family may retain information lost by the minimum in \(v_6\).

For a finite node set \(T\), define a normalized joint-ordering score schematically by
\begin{equation}
  J_H(T)=\sup_{\lambda_{u,v}\geq0}
  \left\{
    \sum_{u,v}\lambda_{u,v}\,F_{(2^u3^v)}(T):
    \sum_{u,v}\lambda_{u,v}(u\log2+v\log3)=1
  \right\},
  \label{eq:ideal-envelope}
\end{equation}
where \(F_{\mathfrak b}(T)\) is the logarithmic generalized factorial at the ideal
\(\mathfrak b\).  The normalization prevents a trivial gain from repeatedly counting the
same local valuation.  Formula \cref{eq:ideal-envelope} is a proposed diagnostic, not a
theorem from the cited papers.

\begin{question}
Does the ideal envelope for the orbit grid exceed the convex envelope of the separate
\((2)\)- and \((3)\)-factorials by a positive leading-order amount?  Equivalently, is there a
nonlinear synchronization phenomenon, or do arbitrary-ideal orderings merely repackage the
same two local valuations?
\end{question}

\subsection{Matroid and greedoid formulations of simultaneous defect}

Grinberg--Petrov show that finite optimality structures inspired by \(p\)-orderings have a
matroidal description, while Krachun--Mirgalimova establish a strong greedoid structure for
removed orderings \cite{GrinbergPetrov,KrachunMirgalimova}.  Johnson gives algorithms for
\(r\)-removed and order-\(h\) variants \cite{Johnson2010}.  These results suggest replacing
the yes/no question ``is there a simultaneous ordering?'' by an optimization problem on
common sets.

For a finite ambient set \(T\subset\Z\) and \(B\subset T\), set
\[
  c_p(B)=v_p\!\left(\prod_{\substack{x,y\in B\\x<y}}(x-y)\right),
  \qquad
  c_p^\star(k)=\min_{\substack{B\subset T\\|B|=k}}c_p(B).
\]
The simultaneous defect of a \(k\)-set is
\[
  \delta_k(B)
  =\bigl(c_2(B)-c_2^\star(k)\bigr)\log2
   +\bigl(c_3(B)-c_3^\star(k)\bigr)\log3,
\]
and the best compromise is
\begin{equation}
  \delta_k^\star=\min_{|B|=k}\delta_k(B).
  \label{eq:matroid-defect}
\end{equation}
The baseline no-simultaneous-ordering result says that \(\delta_k^\star\) is eventually
nonzero under an ordering formulation.  It does not quantify \(\delta_k^\star\), and it does
not exclude common optimal \emph{sets} at selected cardinalities.

A matroid-intersection algorithm can compute \cref{eq:matroid-defect} exactly for finite
truncations, while exchange inequalities may yield asymptotic lower bounds.  There are two
possible payoffs:

\begin{itemize}
\item If \(\delta_k^\star\) is small, the result is negative but clarifying: local conflict
cannot supply the missing determinant scale.
\item If \(\delta_k^\star\) grows on the determinant scale, every common node set pays a
quantified synchronization penalty.  Depending on the sign in the chosen determinant
normalization, that penalty can either force extra divisibility or prove that a proposed
row-selection strategy is impossible.
\end{itemize}

Care is needed with orientation: \(p\)-orderings minimize certain difference valuations,
whereas determinant amplification seeks large valuations after structured row/column
operations.  The computational experiment must trace exactly how the factorial enters the
denominator or determinant bound rather than infer the sign from terminology.

\subsection{Adelic orderings and regular bases}

Chabert--Peruginelli develop adelic versions of ordering and regular-basis theory for rings of
integer-valued polynomials \cite{ChabertPeruginelli}.  Conceptually this is attractive because
one chooses the global basis before decomposing into local factors.  Let \(E_p\) be the
closure of the candidate set in \(\Z_p\), and let \(\underline E=\prod_pE_p\) with suitable
restricted-product conditions.  A successful application would provide a regular basis for
\(\operatorname{Int}(\underline E,\widehat\Z)\) whose leading coefficients encode all places
simultaneously.

The applicability checklist is concrete:

\begin{enumerate}
\item Determine the closures of \(\langle2,3\rangle\) and \(\langle2,M\rangle\) in
\(\Z_p\), including exceptional behavior at \(2\) and \(3\).
\item Verify the finiteness/residual-cardinality hypotheses required for a global regular
basis.  It is not enough to say that the subgroup generated by \(2,3\bmod p\) is usually
large; the exact condition in the theorem must be proved uniformly in degree.
\item Compute the first basis denominators and compare them with the product of ordinary
Bhargava factorials.
\item Determine whether the basis is sensitive to the synchronized exponent pair or only to
the adelic closure, which may forget the orbit ordering completely.
\end{enumerate}

If the closure already factors as an unrestricted product of local closures, no new coupling
will appear.  Conversely, a global degree obstruction or a nonfactorizable leading
coefficient would be precisely the kind of B9 invariant sought by the baseline program.

\subsection{Equidistribution and optimal-ordering asymptotics}

Fr\k{a}czyk--Szumowicz and Byszewski--Fr\k{a}czyk--Szumowicz study simultaneous
\(p\)-orderings, volume minimization, and optimal equidistribution in number fields
\cite{FraczykSzumowicz,ByszewskiFraczykSzumowicz}.  These results help answer a question that
finite computations alone cannot: what limiting measure do near-optimal node sets approach?
If the \(2\)-optimal and \(3\)-optimal measures are distinct, the minimum common defect may
be approximated by an energy-transport cost between them.  This creates a bridge to the
capacity methods of \cref{sec:capacity}.

The useful output would be a variational formula of the form
\[
  \lim_{k\to\infty}\frac{\delta_k^\star}{s(k)}
  =\inf_{\mu}\bigl(I_2(\mu)-I_2^\star+I_3(\mu)-I_3^\star\bigr),
\]
for the correct scale \(s(k)\).  A strictly positive infimum would establish an asymptotic
synchronization cost; a zero infimum would close this family of approaches.

\clearpage
\section{From local divisibility to a global auxiliary relation}
\label{sec:determinant}

\subsection{The classical determinant-method pattern}

Bombieri--Pila begin with points on a short analytic arc and an evaluation determinant of
monomials.  Analytic smoothness makes the determinant small; integrality then forces it to
vanish, placing the points on a low-degree algebraic curve
\cite{BombieriPila,PilaPostulation}.  Heath-Brown replaces part of the real smallness by
\(p\)-adic divisibility, and Salberger globalizes the construction; Walsh sharpens the
resulting curve bounds in a way that removes a logarithmic loss in that algebraic setting
\cite{HeathBrown2002,Salberger2007,Walsh2015}.

The reusable pattern is
\[
  \text{many congruent points}
  \Longrightarrow
  \text{large local determinant valuation}
  \Longrightarrow
  \text{global auxiliary hypersurface}.
\]
It is stronger than proving a valuation lower bound for each determinant term.  The
hypersurface is selected after local clustering and linear algebra, so cancellation is a
feature rather than an uncontrolled remainder.

For a rational evaluation determinant \(D\neq0\), the product formula gives, after clearing
known denominators,
\[
  \log|D|_\infty=\sum_p v_p(D)\log p.
\]
Thus a sufficient local--global inequality is
\begin{equation}
  \sum_{p\in\mathcal P}v_p(D)\log p
  > \log|D|_\infty+\log Q_D,
  \label{eq:product-determinant}
\end{equation}
where \(Q_D\) records any cleared denominator not already included in the local valuations.
Then \(D=0\).  Priority~A asks for precisely such an inequality with
\(\mathcal P\) containing at least \(2\) or \(3\), on a determinant with only
\(O(\log X)\) strategically selected rows.

\subsection{Why the published theorems do not directly apply}

The curve \(\mathcal A=\{(2^t,3^t):0\leq t<1\}\) is analytic and definable but not algebraic.
The full grid \(S_K\) is an integral point set, but it is not given as the bounded-height
rational points on an algebraic variety of fixed degree.  The determinant-method literature
uses the algebraic degree at two critical places:

\begin{enumerate}
\item to ensure that an auxiliary hypersurface not containing the original variety cuts the
point set in controlled degree; and
\item to iterate the covering argument without the auxiliary equations becoming arbitrary.
\end{enumerate}

For the Alaoglu--Erd\H{o}s problem the desired conclusion is itself an auxiliary algebraic
relation on a transcendental orbit.  Consequently, the first vanishing determinant is not
the end of the argument.  One must either obtain a relation with fewer monomials than the
baseline sparse-interpolation lower bound, or obtain enough compatible relations over
several blocks to force an identity.

\subsection{A semigroup-adapted global determinant prototype}

Let
\[
  P_{n,k}=(2^nM^k,3^nA^k)
\]
and choose a finite exponent set \(\Lambda\subset\N_0^2\).  For rows
\(R=\{(n_i,k_i)\}_{i=1}^N\), set
\[
  D(R,\Lambda)=\det\!\left[
    (2^{n_i}M^{k_i})^a(3^{n_i}A^{k_i})^b
  \right]_{i,(a,b)\in\Lambda}.
\]
The baseline tropical lower bound uses
\[
  v_p(D)\geq
  \min_{\sigma}\sum_i v_p\bigl(E_{i,\sigma(i)}\bigr),
\]
which cannot see cancellation between minimizing permutations.  The global determinant
method suggests a different sequence of operations:

\begin{enumerate}
\item Partition exponent pairs \((n,k)\) into residue cells on which the unit parts of
\(2^nM^k\) and \(3^nA^k\) admit \(p\)-adic analytic expansions.
\item Within a cell, replace rows by multivariate finite differences in \(n\) and \(k\).
This exposes powers of \(p\) coming from congruence and derivative order rather than from a
common monomial factor.
\item Choose the monomial lower set \(\Lambda\) to match the finite-difference filtration, as
in a Hilbert-function determinant method.
\item Use the resulting local rank filtration to select one global minor with large total
valuation; only then apply the archimedean bound.
\end{enumerate}

At primes \(p\nmid6MA\), one can restrict \(n,k\) to residue classes modulo the torsion of
\(\Z_p^\times\) and write unit powers through the \(p\)-adic logarithm.  At \(p=2,3\), the
nonunit valuations are explicit linear functions of \((n,k)\), while the residual unit parts
may be handled after a deeper congruence restriction.  This division into valuation and unit
coordinates is standard in \(p\)-adic analytic parameterization and should be made explicit
before determinant searches are run.

\begin{criterion}
A successful prototype must produce, for a family with \(N\to\infty\),
\[
  \sum_{p\in\{2,3\}}v_p(D)\log p-\log|D|_\infty
  \geq c\,s(N)
\]
with \(c>0\) and \(s(N)\) at the leading scale of both terms, after all denominators and
monomial normalizations are included.  A gain caused solely by a removable common row or
column factor is not counted.
\end{criterion}

\subsection{How a vanishing minor could become a proof}

Suppose \(D(R,\Lambda)=0\).  Then a nonzero polynomial supported on \(\Lambda\) vanishes on
\(R\).  There are three possible ways to upgrade this finite relation.

\begin{enumerate}
\item \textbf{Boundary propagation.}  Select overlapping triangular blocks so that the
nullspaces agree on their intersections.  The boundary-dilation equations in the baseline
report are designed for this purpose.
\item \textbf{Shift covariance.}  The transformations
\((n,k)\mapsto(n+1,k)\) and \((n,k+1)\) multiply coordinates by fixed diagonal factors.
A minimal-support nullvector that persists under both shifts may force a Laurent-polynomial
identity on the entire semigroup.
\item \textbf{Fewnomial contradiction.}  Choose \(|\Lambda|\) below the number of distinct
points in a way that violates the baseline monomial lower bound, using extra local
divisibility rather than generic interpolation.
\end{enumerate}

The second route can be formulated in module language.  Let \(I_R\) be the ideal of
polynomials vanishing on a block \(R\).  One seeks a low-support element in an intersection
of shifted ideals
\[
  I_R\cap T_n^{-1}I_R\cap T_k^{-1}I_R,
\]
where \(T_n,T_k\) are the two diagonal semigroup actions.  This makes the passage from a
local determinant to a global relation a finite compatibility problem rather than a vague
compactness argument.

\subsection{Walsh's missing logarithm: analogy and limit}

Walsh removes a logarithmic factor from a bounded-rational-point estimate on algebraic
curves \cite{Walsh2015}.  Since the baseline report also misses by one logarithm, the analogy
is tempting.  The mechanisms are not interchangeable.  Walsh exploits the Hilbert function
and degree of a fixed algebraic curve to choose a better auxiliary polynomial and count its
intersections.  In the present problem there is no fixed algebraic curve of bounded degree;
producing one is the objective.

The transferable question is narrower: does the baseline semigroup determinant use a
rectangular monomial family where a Hilbert-function-adapted staircase would save a boundary
layer?  This is a concrete optimization problem.  It should be tested simultaneously with
the multivariate fixed divisors of \cref{sec:factorials}, because a staircase that is
archimedean-optimal may be locally poor, and vice versa.

\subsection{Habsieger's alternants}

Habsieger's paper is directly concerned with alternants around the four-exponentials
conjecture \cite{Habsieger2007}.  Its value here is not a new unconditional rank theorem but
a mature normal form for the determinant obstruction.  The current semigroup matrices should
be translated into that language for two reasons:

\begin{itemize}
\item to check whether positivity, ordering, or total-positivity hypotheses proved in the
baseline report imply a special alternant inequality not available in the general setting;
\item to avoid rediscovering a determinant transformation whose failure is already encoded
in Habsieger's examples.
\end{itemize}

A useful deliverable is a lemma-by-lemma dictionary between Habsieger's alternants and the
baseline matrices, ending with the first hypothesis that cannot be verified.  That point is
likely to be a sharper formulation of the missing arithmetic input.

\subsection{Recent global-field refinements}

Recent work of Binyamini--Cluckers--Kato and related authors sharpens determinant bounds
uniformly over global fields and improves dimension-growth estimates
\cite{BinyaminiCluckersKato}.  These results reinforce two design choices:

\begin{enumerate}
\item keep all places in one global height calculation rather than bolt a \(p\)-adic estimate
onto an archimedean determinant at the end; and
\item choose the monomial/Hilbert-function budget before estimating local determinants.
\end{enumerate}

They do not remove the transcendental-curve mismatch.  Their role is architectural and
quantitative once an algebraic auxiliary family has been manufactured.

\clearpage
\section{Conjugate control, house bounds, and adelic capacity}
\label{sec:capacity}

\subsection{The Newton--Kummer input}

The baseline report constructs shifted higher Newton--Kummer cocycles
\(\Xi_{N,m}\) that are algebraic integers.  In a distinguished real embedding, an adaptive
choice \(N=\ell\), \(m=\ell-1\) gives
\[
  \log|\Xi_\ell|=-c_0\ell+o(\ell),
  \qquad c_0=0.905325\ldots,
\]
while genuine integer exponents make the cocycle exactly zero.  This is unusually strong
archimedean information: a nonzero algebraic integer is exponentially small in a natural
parameter.  The obstruction is that the ambient radical field has degree
\(\ell^\rho\), \(\rho\in\{3,4\}\), and the other conjugates can be
\(\exp(O(\ell))\).  The norm lower bound is therefore far too weak, and ordinary character
projections erase the cross-base cancellation.

The literature in this section is relevant only if it reduces the effective conjugate
complexity.  No house theorem can contradict one small embedding while unconstrained
conjugates are allowed to be enormous.

\subsection{Dimitrov's Schinzel--Zassenhaus theorem}

Dimitrov proves a sharp universal lower bound of Schinzel--Zassenhaus type
\cite{DimitrovSZ}.  In a polynomial formulation, if \(P\in\Z[z]\) has degree \(n\),
\(P(0)=1\), and is not cyclotomic, then a root lies in
\[
  |z|\leq 2^{-1/(4n)}.
\]
Applied to the reciprocal minimal polynomial of a non-root-of-unity algebraic integer
\(\alpha\), this yields
\[
  \overline{|\alpha|}\geq 2^{1/(4n)},
\]
where the left side is the house.

For \(\Xi_\ell\), this theorem is not directly contradictory: the known conjugate upper
bound is exponentially larger than \(2^{1/(4\ell^\rho)}\).  The theorem becomes useful only
after one of the following reductions.

\begin{enumerate}
\item \textbf{Reciprocal normalization with uniform annular control.}
Construct an algebraic integer \(\eta_\ell\) from \(\Xi_\ell\) whose conjugates all lie in
\(e^{o(1)}\leq |z|\leq e^{o(1)}\), while the distinguished conjugate remains
\(e^{-c\ell}\).  Then a reciprocal/house theorem can force a contradiction.

\item \textbf{Low-degree quotient.}
Show that \(\Xi_\ell\), or a normalized ratio of two shifted cocycles, lies in a subfield of
degree \(O(\ell)\) rather than \(\ell^\rho\).  The linear exponential smallness could then
compete with a degree-sensitive lower bound.

\item \textbf{Controlled factor of the conjugate polynomial.}
Prove that the minimal polynomial has a noncyclotomic factor of low degree all of whose roots
inherit the archimedean cancellation.  Character projections do not provide this factor, so
a nonlinear operation is required.
\end{enumerate}

Pritsker's work on algebraic integers with conjugate symmetry and related house bounds may be
closer to a normalized cocycle if a unit-circle symmetry can be forced
\cite{PritskerHouse}.  Again, symmetry without an upper bound is insufficient: reciprocal
roots may simply place one conjugate far outside the unit circle.

\subsection{A conjugate-control ledger}

Every proposed modification of the cocycle should be evaluated by the following ledger.

\begin{center}
\small
\begin{tabularx}{0.96\textwidth}{L{0.23\textwidth}Y}
\toprule
Quantity & Required bound\\
\midrule
Degree & Preferably \(O(\ell)\), or at worst a scale for which the lower bound competes with
\(e^{-c_0\ell}\).\\
Distinguished embedding & Preserve \(\log|\eta_\ell|\leq-c\ell+o(\ell)\) with \(c>0\).\\
Other archimedean embeddings & Uniform \(o(\ell)\), ideally \(O(1)\), after normalization.\\
Nonarchimedean denominators & None for an algebraic integer; otherwise fully included in the
height/product formula.\\
Cyclotomic/unit exceptions & Excluded by a structural argument, not by numerical evidence.\\
Dependence on \(M,A\) & Must retain both outputs; a projection that leaves only one axis has
lost the essential cancellation.\\
\bottomrule
\end{tabularx}
\end{center}

This ledger prevents a common false gain: dividing by a large algebraic factor improves the
distinguished embedding but introduces a denominator of the same global height.

\subsection{Transfinite diameter and resultant identities}

DeMarco--Rumely give transformation laws for transfinite diameter under polynomial maps,
with the resultant recording the arithmetic distortion \cite{DeMarcoRumely}.  Rumely's
adelic capacity and Fekete--Szeg\H{o} theorems with local rationality conditions provide a
global framework for sets of possible conjugates \cite{RumelyFeketeSzego}.  These results
can be used in two complementary ways.

\subsubsection{Positive use: forbid a conjugate configuration}

Associate to each place \(v\) a compact set \(E_v\) containing every conjugate of a
normalized cocycle.  At the distinguished real place, \(E_\infty\) contains a tiny disk; at
other places, integrality may give unit disks or explicitly enlarged disks.  If the global
capacity
\[
  \gamma(\mathbb E)=\prod_v\gamma_v(E_v)
\]
is strictly less than one, only finitely many algebraic numbers of the allowed type can have
all conjugates in the adelic set.  A family of growing degree would then be impossible unless
it falls into a special algebraic locus.

The hard step is making the small real disk influence enough conjugates.  One tiny component
among \(\ell^\rho\) archimedean embeddings contributes too little.  A successful construction
must cause a positive proportion of embeddings to lie in controlled small sets, or reduce
the degree before capacity is applied.

\subsubsection{Negative use: certify capacity criticality}

The same calculation can show that an entire determinant/cocycle strategy is incapable of
winning.  If the optimal adelic sets have global capacity exactly one, the product formula is
at equilibrium; no rearrangement of the same local data can produce a strict contradiction.
This would turn the baseline's observed one-logarithm deficit into a theorem about the
capacity of the orbit configuration.

A recommended computation is therefore:

\begin{enumerate}
\item derive the limiting archimedean logarithmic energy of the normalized cocycle or
Vandermonde node measure;
\item compute the \(2\)- and \(3\)-adic capacities of the corresponding closures;
\item include the resultant/discriminant factors exactly; and
\item test whether the global product is \(<1\), \(=1\), or \(>1\).
\end{enumerate}

This is also a disciplined way to integrate the equilibrium-measure calculations already
present in the baseline report with established potential theory.

\subsection{Capacity and generalized-ordering limits}

The generalized factorials of \cref{sec:factorials} and transfinite diameter are two finite
and asymptotic faces of the same phenomenon: products of pairwise differences determine both
Vandermonde valuations and logarithmic energy.  If the optimal-ordering measures at \(2\)
and \(3\) converge, one can ask whether a single global equilibrium measure minimizes
\[
  I_\infty(\mu)+I_2(\mu)+I_3(\mu)
\]
subject to the semigroup support constraints.  A strict gap between the minimum over all
measures and the minimum over orbit-realizable measures would be a genuine local--global
obstruction.  No cited paper supplies this semigroup-constrained variational theorem, but the
combined literature suggests exactly how to formulate it.

\begin{stopbox}[Stop condition for house/capacity methods]
If every normalized cocycle retains degree \(\ell^\rho\) with only one controlled small
embedding, and the adelic capacity calculation is at least one, then Schinzel--Zassenhaus
bounds cannot see the construction.  Further improvements to the numerical house constant
would not change that conclusion.
\end{stopbox}

\section{Parametric geometry, Pad\'e systems, and potential theory}
\label{sec:pade}

\subsection{Parametric geometry as a language for the bilinear wall}

Schmidt--Summerer parametric geometry of numbers and Roy's realization/classification of
\(n\)-systems encode Diophantine approximation through the successive minima of a
one-parameter family of convex bodies \cite{SchmidtSummerer,RoyPGN}.  The modern variational
formulation of Das--Fishman--Simmons--Urba\'nski makes spectra and extremal trajectories
systematic \cite{DFSU2024}.

For the Alaoglu--Erd\H{o}s problem, one can form a lattice or pair of lattices encoding
simultaneously
\[
  m\log2-n\log3
  \qquad\text{and}\qquad
  a\log M-b\log A.
\]
The exact relation \(\log M\log3=\log A\log2\) couples the slopes of these two approximation
problems.  A PGN reformulation is useful only if integrality forces an impossible pattern in
the combined successive minima.  Merely restating that each logarithmic ratio has good
continued-fraction approximants does not add information.

A concrete model is to introduce a four-dimensional lattice generated by rows containing
\[
  (1,0,\log2,\log M),
  \qquad
  (0,1,\log3,\log A),
\]
with anisotropic weights chosen so that short vectors correspond to simultaneous small
linear forms.  The terminal identity imposes a rank-one condition on the real \(2\times2\)
logarithm matrix, while unique factorization and primitive-output constraints impose
integral restrictions on admissible lattice vectors.  The desired theorem would say that no
self-similar \(n\)-system realizes all those restrictions.

\begin{warningbox}[What PGN does not provide automatically]
Parametric geometry classifies possible approximation exponents for arbitrary real vectors.
The present vector is special because its exponentials are integers.  Unless that arithmetic
condition is translated into a restriction on the minima trajectory, PGN is only a change of
coordinates for the same open four-exponentials case.
\end{warningbox}

\subsection{Shifted logarithms and Pad\'e constructions}

Kawashima--Po\"els and recent work on shifted logarithms develop explicit Pad\'e systems and
linear-independence measures for families of logarithms
\cite{KawashimaPoels,HirataKohnoMuroiWashio}.  Related \(S\)-unit Pad\'e constructions aim to
control denominators and linear forms simultaneously \cite{SUnitPade}.  These methods can
contribute in two restricted ways:

\begin{enumerate}
\item improve effective lower bounds for combinations of \(\log2\) and \(\log3\), thereby
extending finite verified regions; or
\item supply an auxiliary table whose denominator structure can be synchronized with the
candidate semigroup, creating a second integer rather than merely a better real
approximation.
\end{enumerate}

The first payoff is Priority~H and is realistic.  The second would be genuinely new.  To test
it, choose Pad\'e approximants to a vector of functions containing
\[
  \log(1-z),\quad \log(1-2z),\quad \log(1-3z)
\]
or a binomial family, specialize at rational points whose denominators are powers of \(2\)
and \(3\), and track the exact denominator ideal.  The specialization must produce a relation
in which the unknown \(M,A\) enter algebraically; otherwise it reduces to a known estimate
for \(\log2/\log3\).

\subsection{Hankel determinants and irrationality criteria}

Bugeaud--Han--Wen--Yao and related work connect Hankel determinants, Pad\'e approximation,
and irrationality of power-series values \cite{BugeaudHanWenYao}.  The baseline report
already studies Hankel/Loewner and Markov--Pad\'e defects.  The literature can sharpen that
program by separating three questions:

\begin{enumerate}
\item Is the moment sequence nondegenerate, so that infinitely many Hankel determinants are
nonzero?
\item What are the exact denominator valuations of those determinants at \(2\) and \(3\)?
\item Does potential theory give an asymptotic constant small enough to contradict
integrality after the denominator is cleared?
\end{enumerate}

Nonvanishing alone is not enough.  A useful Hankel determinant must have an arithmetic
normalization whose denominator cost is smaller than its archimedean decay.  This is the same
local--global ledger as \cref{eq:product-determinant}.

\subsection{Classical Pad\'e potential theory}

The baseline report derives logarithmic energies for projected lattice triangles and for a
Stieltjes-type interpolant.  Gonchar--Rakhmanov equilibrium measures and the Stahl--Totik
theory of rational approximation provide established tools for identifying exact nth-root
asymptotics \cite{GoncharRakhmanov,StahlTotik}.  Beukers's Pad\'e constructions for binomial
functions offer explicit comparison families \cite{BeukersPade}.

This literature can answer a crucial diagnostic question: is the residual positive constant
in the current Pad\'e calculation an artifact of the chosen nodes, or is it the equilibrium
energy of the relevant compact set and therefore asymptotically optimal?  The workflow is:

\begin{enumerate}
\item identify the compact set and external field determined by the node distribution;
\item solve or cite the equilibrium problem;
\item derive the asymptotic zero distribution of the Pad\'e denominator;
\item compute the exact exponential rate of the remainder; and
\item compare it with the arithmetic denominator rate.
\end{enumerate}

If the existing construction already attains the equilibrium rate and still loses, no
reweighting of the same support can help.  A successful improvement would then have to
change the arithmetic object or add a second orthogonality condition, not merely tune
coefficients.

\clearpage
\section{Denominator-sensitive counting and unlikely intersections}
\label{sec:counting}

\subsection{The quantitative mismatch with ordinary counting}

The compact orbit
\[
  (U_k,V_k)=\bigl(2^{\{k\beta\}},3^{\{k\beta\}}\bigr)
\]
contains \(\asymp\log B\) points of multiplicative height at most \(B\), because logarithmic
height grows linearly with \(k\).  A Pila--Wilkie estimate of the form \(O(B^\eps)\) therefore
does not constrain the orbit at all: \(\log B\ll B^\eps\).  Effectivizing the constant in the
same estimate does not change this comparison.

The arithmetic information omitted by ordinary height is unusually rigid.  Up to reductions
controlled by fixed valuations of \(M\) and \(A\), the first coordinate has denominator
supported only at \(2\), the second only at \(3\), and both denominator exponents are governed
by the same floor \(\lfloor k\beta\rfloor\).  A theorem capable of proving the conjecture
would resemble
\begin{equation}
  \#\left\{
    t\in[0,1):
    2^t\in\Z[1/2],\ 3^t\in\Z[1/3],\
    h(2^t,3^t)\leq H,
    \text{synchronized denominators}
  \right\}=o(H),
  \label{eq:support-count}
\end{equation}
when \(H\) is logarithmic height.  The hypothetical orbit gives \(\asymp H\) points.  No
survey located a theorem with this support-and-synchronization sensitivity.

\subsection{Arithmetic unlikely intersections in powers of \texorpdfstring{\(\Gm\)}{Gm}}

Campagna--Dill--Wilms formulate and prove cases of arithmetic unlikely-intersection bounds in
powers of the multiplicative group, including a complete dimension-one result and partial
dimension-two results \cite{CampagnaDillWilms}.  Their framework is closer to the desired
arithmetic packaging than generic o-minimal counting because subgroup complexity and
integral models enter the bound.

There are nevertheless two mismatches.

\begin{enumerate}
\item The point set in the present problem lies on a transcendental analytic arc, not an
algebraic subvariety of \(\Gm^2\).
\item The key restriction is not merely membership in a finite-rank subgroup.  It is the
coordinatewise prime support together with a common exponent parameter.
\end{enumerate}

A viable adaptation would first algebraize a finite orbit block by a low-degree auxiliary
relation obtained from \cref{sec:determinant}, then apply an arithmetic
unlikely-intersection bound to that auxiliary curve and the relevant subgroup schemes.  This
is a two-stage program: the new paper may strengthen the second stage, but it does not
provide the first.

\subsection{Effective Pfaffian counting}

Binyamini--Jones--Schmidt--Thomas give effective versions of Pila--Wilkie-type counting for
Pfaffian structures \cite{BinyaminiJonesSchmidtThomas}.  The arc \(t\mapsto(2^t,3^t)\) is
well within such tame analytic settings.  The result is useful for explicit bookkeeping once
one has a counting exponent below the orbit density.  By itself, it retains only height and
definable complexity, so the bound remains too large.

The most plausible way to combine it with arithmetic is to partition rational points by
exact denominator ideals and apply determinant or congruence arguments inside each fiber.
However, there are only \(O(H)\) relevant synchronized fibers up to logarithmic height \(H\),
and one orbit point per fiber is already enough.  A successful theorem must show that most
fibers are empty for arithmetic reasons; a uniform bound per fiber is not sufficient.

\subsection{Integer-valued definable functions}

Jones--Thomas--Wilkie and later work on integer-valued o-minimal functions establish strong
polynomiality or growth dichotomies for definable functions taking integer values on the
integers \cite{JonesThomasWilkie,BhardwajEtAl}.  This literature offers a clean conditional
route:

\begin{quote}
Construct a function \(F\) definable in the required o-minimal structure, with controlled
growth, such that \(F(n)\in\Z\) for every \(n\in\N\) and such that \(F\) encodes
\(n\mapsto n^\theta\) or the candidate semigroup.  The definable integer-valued theorem may
then force \(F\) to be a polynomial, which is incompatible with irrational \(\theta\).
\end{quote}

The current canonical interpolants are selected through discrete minimization or infinite
products on sparse nodes and are not known to be definable in a useful uniform structure.
Thus the missing lemma is exactly the construction of a definable representative.  The
literature should not be applied to \(x^\theta\) merely because that function itself is
definable: it is not integer-valued on every integer, only on the hypothetical candidate
submonoid.

A more promising definable object may be an envelope or extremal function whose parameter
set is continuous and whose integrality on \(\Cand\) forces integrality on a definable family.
This would convert sparse interpolation into a unary integer-valued problem, but no such
extension principle is currently known.

\subsection{Analytic gcd bounds}

Pasten--Wang extend gcd estimates to analytic/meromorphic functions using Nevanlinna theory
\cite{PastenWang}.  The baseline program contains gcd/lcm lattices, common-divisor heuristics,
and pairs of interpolating functions.  The theorem could become relevant if one constructs
algebraically independent meromorphic functions \(F,G\) for which the counterexample forces
large common divisors or many common zeros at the candidate nodes.

The required translation should be written before the theorem is invoked:

\begin{enumerate}
\item identify \(F\) and \(G\) and prove the algebraic-independence hypothesis;
\item express the integer gcd or common-zero multiplicity in the Nevanlinna counting
function used by the theorem;
\item show that the candidate orbit forces a lower bound exceeding the theorem's upper bound;
and
\item include poles and denominator contributions exactly.
\end{enumerate}

Finite differences of one interpolant are not automatically algebraically independent from
it, and sharing a discrete set of integer values is not the same as sharing zeros.  This is a
conditional attachment point, not a direct gcd theorem for \(M^n-1,A^n-1\).

\subsection{Why generic BMZ, Maurin, and Zilber--Pink remain out of scope}

Bombieri--Masser--Zannier and Maurin prove finiteness or structural results for intersections
of an \emph{algebraic} curve with algebraic subgroups of \(\Gm^n\)
\cite{BMZ1999,Maurin2008}.  The map \(t\mapsto(2^t,3^t)\) does not define such a curve.
Embedding its graph in a larger exponential-algebraic structure does not restore the
algebraic degree needed by the theorem.  These results become relevant only after an
auxiliary algebraic curve has been forced, in which case they may help show that the curve
cannot contain too many orbit points unless it is a subgroup translate.

\begin{stopbox}[Stop condition for counting]
Do not pursue a counting theorem whose bound depends only on ordinary height and definable
complexity.  To beat the hypothetical orbit it must use denominator support, synchronized
exponents, subgroup complexity after algebraization, or another arithmetic invariant that
makes the right side of \cref{eq:support-count} sublinear in logarithmic height.
\end{stopbox}

\section{Complex powers and model-theoretic existence results}
\label{sec:powers}

Gallinaro proves exponential-algebraic closedness results for systems of raising-to-powers
type, and Gallinaro--Kirby establish quasiminimality phenomena for structures with complex
power maps; Aslanyan--Kirby--Mantova give related geometric existence theorems for exponential
equations \cite{Gallinaro2025,GallinaroKirby,AslanyanKirbyMantova}.  These papers are
relevant because the distinguished power map
\[
  z\longmapsto z^\theta,
  \qquad \theta=\frac{\log3}{\log2},
\]
is exponentially algebraic even though \(\theta\) is transcendental.

Their conclusions point in the opposite direction from the desired theorem.  They show that
suitable systems have solutions and that the ambient structure is tame; they do not show
that a prescribed solution with integer coordinates is impossible.  In particular, an
existence theorem for auxiliary powers does not manufacture a multiplicatively independent
integer base \(q\) with controlled \(q^x\).

There are still two diagnostic uses.

\begin{enumerate}
\item Formulate the counterexample locus in the language of complex powers and determine its
predimension.  If the integer-support constraints force an atypical predimension drop beyond
what the existence theory permits generically, this may identify the exact additional
arithmetic statement required.
\item Use quasiminimality/tameness to rule out proposed definability arguments that would
make the set of candidates too complicated.  This can prevent a false model-theoretic route,
but is unlikely to yield a quantitative contradiction.
\end{enumerate}

The appropriate classification is therefore \badgeTemplate/\badgeDiagnostic, not a proposed
source of the missing independent column.

\section{Function-field, Mahler, and digit/carry analogues}
\label{sec:analogues}

\subsection{Positive-characteristic logarithm independence}

Papanikolas's Tannakian theory for Anderson--Drinfeld motives identifies transcendence degree
with the dimension of a difference-Galois group, and Chang--Papanikolas determine algebraic
relations among periods and logarithms of Drinfeld modules
\cite{Papanikolas2008,ChangPapanikolas}.  Recent work of Gezmi\c{s}--Namoijam extends parts of
this picture to logarithms of Anderson \(t\)-modules \cite{GezmisNamoijam}.  In positive
characteristic these are unconditional analogues of transcendence statements that remain
open over \(\Q\).

The mechanism replacing the missing archimedean size is not a stronger determinant estimate.
It is Frobenius-semilinear difference structure: a rigid analytic trivialization satisfies a
\(\sigma\)-difference equation with algebraic coefficients, and the dimension of the
associated Galois group controls algebraic independence.  This observation suggests a
structural design criterion for characteristic zero:

\begin{quote}
A successful auxiliary object should possess a semilinear or differential symmetry whose
Galois group remembers both logarithmic directions and whose arithmeticity is independent of
the hypothetical transcendental parameter.
\end{quote}

The maps \(n\mapsto2n\) and \(n\mapsto3n\), or the two shifts on the exponent grid, are
obvious candidates for a difference structure.  They do not behave like Frobenius: there is
no coefficient-field endomorphism making \(2^z\) and \(3^z\) solutions of one common
semilinear equation.  A transfer from the function-field proof would therefore require a new
charateristic-zero difference module, not an analogy at the level of statements.

A useful negative deliverable is to prove that the natural two-shift module generated by the
candidate table has difference-Galois group of the expected rank unless
\(\beta\in\Q\).  Such a theorem might repackage the conjecture but would at least reveal
whether the missing condition is a Galois-group dimension statement.

\subsection{Mahler's method and automatic structure}

Adamczewski--Faverjon give powerful theorems on linear and algebraic relations among values of
Mahler functions at algebraic points, with effective variants
\cite{AdamczewskiFaverjon2017,AdamczewskiFaverjon2018}.  Adamczewski explicitly discusses the
Alaoglu--Erd\H{o}s problem in the surrounding Mahler/Cobham circle
\cite{AdamczewskiMahlerProblem}.

This literature becomes actionable only after deriving a functional equation such as
\[
  \mathbf F(z)=A(z)\mathbf F(z^q)
\]
with \(A(z)\) algebraic/rational and an algebraic specialization.  The natural candidate
sequences---digits of \(\lfloor k\beta\rfloor\), carries in the \(2\)- and \(3\)-adic
normalizations, or coefficients of an interpolant---do not automatically satisfy a Mahler
equation.  Irrational rotations are Sturmian rather than automatic in general.

There are three precise tests:

\begin{enumerate}
\item Does simultaneous integrality force the carry word to be \(2\)-automatic or
\(3\)-automatic?
\item Does the same sequence admit compatible descriptions in both bases, so that a Cobham-
type theorem forces eventual periodicity?
\item Can a generating function of the orbit be shown to be Mahler without coefficients
containing \(M,A\), or \(\beta\) transcendently?
\end{enumerate}

If the answer to the first two questions is yes, eventual periodicity would likely imply
\(\beta\in\Q\), finishing the problem.  If the sequence remains merely Sturmian, the Mahler
literature is not engaged.

\subsection{Ternary digits of powers of two and carry certificates}

Lagarias studies ternary expansions of powers of two through \(3\)-adic dynamics
\cite{Lagarias2009}.  Saye gives recursive certificates for very large finite ranges of
prescribed trailing ternary patterns, and Dimitrov--Howe analyze powers of \(3\) with few
nonzero binary digits \cite{Saye2022,DimitrovHowe}.  These works do not imply the
Alaoglu--Erd\H{o}s conjecture, but they offer two transferable tools.

First, they show how multiplicative dynamics can be encoded as compact trailing-digit
recursions rather than enormous integer computations.  This is attractive for Priority~I:
a candidate search or zero-free verification can be accompanied by a short recursive
certificate checked in Lean.

Second, the \(3\)-adic Cantor-set perspective suggests studying the set of possible unit
parts of \(M^k\) simultaneously in \(\Z_2\) and \(\Z_3\).  The desired statement would be an
intersection-dimension or entropy bound for an orbit constrained by two incompatible digit
alphabets.  Existing digit-omission theorems treat one base action at a time; they do not
couple the common floor parameter.

The experimental target should therefore be a \emph{synchronized carry automaton}.  For a
finite depth \(r\), record the state
\[
  \bigl(M^k\bmod 2^r,\ A^k\bmod3^r,\ \lfloor k\beta\rfloor\bmod q\bigr)
\]
with transitions in \(k\).  Search for forbidden state combinations stable as \(r\) grows.
Any discovered invariant must be converted into an exact congruence theorem; entropy or
apparent chaos is not evidence by itself.

\subsection{Recent gap/word results}

Hultberg's recent gap principles on semiabelian varieties and Stephan's work on steering
words are remote but potentially useful diagnostics \cite{Hultberg2026,Stephan2026}.  The
former would require an algebraic semiabelian orbit of the correct type; the latter would
require a finite-word encoding whose low complexity follows from the counterexample.  Neither
input is currently available.  They belong in the lead inventory so that a future structural
breakthrough can activate them, but they should not compete with the direct programs in
\cref{sec:pila,sec:factorials,sec:determinant}.

\begin{takeawaybox}[Template lesson]
The positive-characteristic and Mahler successes all possess a rigid functional equation
over an arithmetic coefficient field.  The Alaoglu--Erd\H{o}s table has many shift identities,
but their coefficients involve the unknown integral outputs.  The decisive transfer problem
is to eliminate that dependence while retaining both logarithmic directions.
\end{takeawaybox}

\section{Rank theorems and the four-exponentials frontier}
\label{sec:rank}

\subsection{Conditional resolution and unconditional boundary}

The Four Exponentials Conjecture says that if \(x_1,x_2\) and \(y_1,y_2\) are each
\(\Q\)-linearly independent, then at least one of \(e^{x_i y_j}\) is transcendental.  Apply
it to
\[
  (x_1,x_2)=(1,x),
  \qquad
  (y_1,y_2)=(\log2,\log3).
\]
If all four exponentials \(2,3,2^x,3^x\) are algebraic, Four Exponentials forces
\(x\in\Q\), and unique factorization then gives \(x\in\Z\).  Thus the conjecture is an exact
special case of the four-exponentials frontier, not merely adjacent to it.

Dasgupta's Part~I gives a clean modern account of unconditional rank theorems for matrices of
logarithms of algebraic numbers, including Baker and Waldschmidt--Masser results and Roy's
structural-rank perspective \cite{DasguptaPartI}.  Dasgupta--Kakde's Matrix Coefficient
Conjecture makes the balanced \(2\times2\) obstruction explicit
\cite{DasguptaKakdePartII}.  These sources are important for calibrating claims: a proposed
rank argument should be checked against the precise hypotheses and known lower bounds before
being presented as new progress.

\subsection{Quadratic independence of prime logarithms}

Equation \cref{eq:terminal-wall} is excluded by a suitable \(\Q\)-linear independence
statement for products \(\log p\log q\).  For example, linear independence of the relevant
part of the symmetric square of the \(\Q\)-space spanned by prime logarithms would force all
coefficients \(m_p,a_p\) to vanish in the prohibited combination.  This is an exact
sufficient condition, but there is no evidence that it is easier than the original problem;
it is a broad fragment of the four-exponentials circle.

The useful research question is narrower: do positivity and disjoint support of the exponent
vectors in \cref{eq:terminal-wall} define a cone on which independence can be proved without
settling arbitrary signed relations?  None of the four surveys found such a theorem.  This
``positive cone'' formulation should be kept visible because determinant, capacity, or
convexity arguments may exploit it even when general quadratic independence remains open.

\subsection{Analytic subgroup and period-matrix packaging}

W\"ustholz's analytic subgroup theorem and quantitative period theorems of
Gaudron--R\'emond provide the natural algebraic-group language for unexpected rank drops
\cite{Wustholz1989,GaudronRemond}.  They are extraordinarily powerful when the logarithms or
periods lie in the Lie algebra of an algebraic group and the relevant coefficients are
algebraic.  In the present matrix, the unknown logarithms \(\log M,\log A\) are logarithms of
algebraic numbers, but the rank-one relation involves the transcendental coefficient ratio
\(\log3/\log2\).  The analytic subgroup theorem does not automatically lower the rank from
two to one.

This packaging may nevertheless interface better with valuation lattices and 1-motives than
raw linear forms do.  A concrete adaptation would construct an algebraic subgroup from the
integral semigroup action, not assume one from the real logarithmic relation.  Until that
subgroup is exhibited, the AST is a frontier theorem rather than a proof route.

\subsection{Roy, Diaz, and the ceiling of current independence results}

Roy's small-value and parametric-geometry results sharpen the six-exponentials frontier, while
Diaz proves substantial algebraic independence for iterated powers with algebraic exponents
\cite{RoySixExp,Diaz1989}.  These theorems show how much independence can be extracted when
an algebraic exponent or a third direction is available.  A hypothetical counterexample is
balanced to remove exactly those extra degrees of freedom.

The correct use of these results is to formulate gateways.  For example:

\begin{itemize}
\item any construction of a third multiplicatively independent integer \(q\) with algebraic
\(q^x\) invokes six exponentials;
\item any construction of a nonzero algebraic \(\gamma\) with algebraic
\(e^{\gamma x}\) invokes five-exponentials-type restrictions; and
\item any reduction of \(\beta\) to an algebraic irrational invokes stronger independence
theorems and likely gives an immediate contradiction.
\end{itemize}

Without such a gateway, quoting the theorem only restates why the two-base case is hard.

\subsection{Triage of the broad rank/dynamics survey}

The useful content of \Rtwo\ is its clean conditional Four Exponentials reduction, its
emphasis on structural rank, and its pointer to the matrix-coefficient framework.  The
following claims are downgraded:

\begin{itemize}
\item Baker--type improvements do not cross the logarithmic-coefficient wall.
\item BCZ/S-unit finiteness does not imply nonexistence and usually assumes a fixed
algebraic curve or recurrence.
\item Symbolic chaos of fractional parts is heuristic unless exact automatic structure is
derived.
\item Generic unlikely intersections and the analytic subgroup theorem require algebraic
input absent from the transcendental arc.
\item Metric B\'ezout or small-value theorems need a new auxiliary construction; they do not
geometrize the problem for free.
\end{itemize}

This correction is not a dismissal of the surrounding theories.  It locates the precise
new lemma each would need before it can touch the conjecture.

\clearpage
\section{Effective finite progress, formal certificates, and the CA interface}
\label{sec:finite}

\subsection{Explicit approximation to \texorpdfstring{\(\log2/\log3\)}{log 2 / log 3}}

Simons and de Weger develop effective continued-fraction and linear-form estimates involving
\(\log2\) and \(\log3\) \cite{SimonsDeWeger,Simons2005}.  These estimates do not address the
bilinear relation \cref{eq:terminal-wall}, but they can make finite exclusion substantially
cheaper.  The baseline verification tests, for each integer \(m\) in a range, whether
\[
  m^\theta,\qquad \theta=\frac{\log3}{\log2},
\]
can be an integer, using rigorous enclosures obtained from convergents and semiconvergents.
A uniform irrationality/linear-form bound can replace many ad hoc precision tiers by one
certificate schema.

The correct benchmark is not merely a larger floating-point scan.  A useful result should
produce a compact proof object of the form
\[
  \forall m\leq X,\qquad
  \operatorname{dist}(m^\theta,\Z)>\varepsilon_m>0,
\]
where every \(\varepsilon_m\) follows from a small list of rational interval enclosures and
integer inequalities.  The certificate must replay in the kernel without trusting MPFR
rounding or an external search result.

Pad\'e and hypergeometric improvements surveyed by Zudilin and developed by Marcovecchio may
reduce the size of these certificates \cite{ZudilinSurvey,Marcovecchio2009}.  Their role
remains finite unless the resulting measure is strong enough to compare with the
exponentially small gap generated by an arbitrary unknown \(M\), which ordinary
linear-logarithm measures are not.

\subsection{Certificate translation and DSLean}

DSLean provides infrastructure for translating certificates from external domain-specific
solvers, including interval arithmetic/Gappa-style proofs, into Lean terms checked by the
kernel \cite{DSLean2026}.  The baseline project already has generated tables and
semiconvergent enclosures.  The natural integration is:

\begin{enumerate}
\item an external generator selects convergents, semiconvergents, and interval splits;
\item Gappa or an exact rational solver verifies each transcendental enclosure from a small
library of analytic lemmas;
\item DSLean translates the certificate into a Lean proof term; and
\item one generic theorem lifts the per-row inequalities to the finite zero-free region.
\end{enumerate}

This architecture separates trusted mathematics from untrusted search while avoiding a
monolithic \texttt{decide} table.  It is also reusable for determinant sign checks, capacity
constants, and Pad\'e remainder inequalities.

\subsection{Digit recursion as proof compression}

Saye's trailing-ternary-digit recursion verifies enormous finite ranges using a short
recursive construction rather than enumerating every exponent \cite{Saye2022}.  A comparable
compression may be possible for the Alaoglu--Erd\H{o}s scan if candidate intervals can be
organized by a common prefix in the binary expansion of \(\theta\log_2m\) or by a common
semiconvergent tier.  The output should be a tree whose internal node carries one interval
inequality proving all descendant values impossible.  This is more promising than simply
raising the numerical frontier because the same tree can be replayed and minimized.

\subsection{The colossally abundant consequence}

The historical motivation should be kept logically separate from the proof search.  Work of
Caveney--Nicolas--Sondow records the consequence that, under the Alaoglu--Erd\H{o}s
assumption, quotients of consecutive colossally abundant numbers are prime
\cite{CaveneyNicolasSondow}.  The original Alaoglu--Erd\H{o}s analysis already restricts such
quotients sharply.

This is an excellent formalization target because it is downstream and combinatorial.  A
clean module should state:

\begin{enumerate}
\item the definition and ordering of colossally abundant numbers;
\item the formula determining their prime exponents from a parameter;
\item the reduction of a composite consecutive quotient to a rational-power equality; and
\item the invocation of the integer-exponent conjecture to force primality.
\end{enumerate}

Formalizing this implication will not advance the transcendence barrier, but it gives a
machine-checked historical payoff and tests whether the conjecture has been stated in the
strongest form needed by the application.

\begin{takeawaybox}[Role of computation]
Finite verification should be optimized for theorem discovery and proof compression.  The
highest-value computations are those that estimate a leading asymptotic constant, expose a
Smith-profile pattern, or produce a reusable exact certificate.  Extending a decimal
frontier without changing the structure remains evidence only.
\end{takeawaybox}

\section{Integrated research program}
\label{sec:program}

The merged literature supports a staged program in which every work package has a binary or
quantitative decision gate.  The packages are ordered to maximize information gained before
large formal or theoretical investments.

\subsection{WP1: Pila growth audit}

\textbf{Goal.}  Decide whether the canonical entire-function program enters Pila's direct
order-two dependence theorem.

\textbf{Tasks.}
\begin{enumerate}
\item Select two canonical functions: one interpolating on \(\Sol\), one on \(\Cand\).
\item Prove maximum-modulus upper bounds using the exact node-counting measure, not a generic
Hadamard product estimate.
\item Obtain a matching lower bound from the concordant envelope or logarithmic potential.
\item Compare the resulting type with \cref{eq:pila-threshold}.
\item If below threshold, classify the algebraic dependence over \(\C(2^z,3^z)\).
\end{enumerate}

\textbf{Decision gate.}  Continue only if the threshold is met, approached within a
structurally improvable factor, or the algebraic-dependence classification yields a new
constraint.  Otherwise record the exact excess constant as a no-go benchmark.

\subsection{WP2: mixed-factorial scaling experiment}

\textbf{Goal.}  Determine whether multivariate/composite/adelic orderings create a leading-
scale synchronization term absent from scalar primewise factorials.

\textbf{Tasks.}
\begin{enumerate}
\item Implement the Prasad--Rajkumar--Reddy and Evrard factorials for lower sets
\(\Lambda\subset\N_0^2\).
\item Compute exact Smith forms for \(S_K\) with symbolic valuation parameters where
possible and concrete primitive pairs otherwise.
\item Implement \(b\)- and ideal-orderings for \(b=2^u3^v\).
\item Solve the simultaneous-set defect \cref{eq:matroid-defect} using matroid/greedoid
algorithms.
\item Fit normalized asymptotics and prove the first stable recurrence or energy formula.
\end{enumerate}

\textbf{Decision gate.}  Continue only if a mixed gain survives after normalization at the
archimedean determinant scale.  A boundary-order gain is insufficient.

\subsection{WP3: global \texorpdfstring{\(p\)}{p}-adic determinant prototype}

\textbf{Goal.}  Turn semigroup congruence into an auxiliary polynomial before applying the
archimedean bound.

\textbf{Tasks.}
\begin{enumerate}
\item Develop \(p\)-adic analytic coordinates for exponent cells at \(p=2,3\) and at one
generic odd prime.
\item Design finite-difference row bases matched to staircase monomial sets.
\item Prove a local rank-filtration lemma analogous to the p-adic determinant method.
\item Combine places globally and calculate every denominator in
\cref{eq:product-determinant}.
\item Test shift covariance of the resulting nullvectors on overlapping blocks.
\end{enumerate}

\textbf{Decision gate.}  The first prototype must beat the tropical termwise ceiling on the
leading scale.  If a general inequality proves that it never can, promote that inequality to
a no-go theorem and stop this family.

\subsection{WP4: CDT holonomy suitability audit}

\textbf{Goal.}  Determine whether arithmetic holonomy can be made to see the
\((\log2,\log3)\) product rather than a neighboring family.

\textbf{Tasks.}
\begin{enumerate}
\item Reproduce the product-of-logarithms theorem with all constants and identify the source
of the near-diagonal restriction.
\item Test telescoping logarithmic decompositions and parameter deformations for uniformity.
\item Search for an arithmetic holonomic generating series of semigroup minors or Pad\'e
defects.
\item Verify coefficient algebraicity and denominator growth before any analytic estimate.
\end{enumerate}

\textbf{Decision gate.}  Continue only if the desired specialization remains inside a
uniform arithmetic family.  Analytic approximation by admissible parameters is not enough.

\subsection{WP5: Newton--Kummer conjugate compression}

\textbf{Goal.}  Reduce the effective degree or control a positive proportion of conjugates of
the exponentially small cocycle.

\textbf{Tasks.}
\begin{enumerate}
\item Search nonlinear resultants, ratios of adjacent shifts, and norm-one normalizations
that retain both \(M\) and \(A\).
\item Track the conjugate-control ledger of \cref{sec:capacity} exactly.
\item Compute local and archimedean capacities of the normalized conjugate sets.
\item Apply Dimitrov/Pritsker only after the degree and annular bounds compete with
\(e^{-c_0\ell}\).
\end{enumerate}

\textbf{Decision gate.}  Stop if only one of \(\ell^\rho\) embeddings is small and global
capacity remains at least one.

\subsection{WP6: certified finite expansion and historical interface}

\textbf{Goal.}  Extend the kernel-replayed region while improving proof architecture.

\textbf{Tasks.}
\begin{enumerate}
\item Integrate explicit continued-fraction bounds with a DSLean/Gappa certificate
translator.
\item Replace flat tables by recursive interval or digit certificates.
\item Use the same infrastructure for numerical constants in WPs 1--5.
\item Formalize the colossally abundant quotient implication in a separate downstream
module.
\end{enumerate}

\textbf{Decision gate.}  A new finite range is worth merging only if its certificate size and
replay cost scale substantially better than row-by-row enumeration.

\subsection{Dependency graph}

\begin{center}
\small
\begin{tabularx}{0.94\textwidth}{L{0.18\textwidth}L{0.28\textwidth}Y}
\toprule
Package & Can start immediately? & Dependencies / shared outputs\\
\midrule
WP1 Pila & Yes & Shares potential-theory constants with WP5 and exact interpolation code with
WP2.\\
WP2 factorials & Yes & Supplies local divisibility data to WP3 and equilibrium measures to
WP5.\\
WP3 determinant & Yes, finite prototype & Benefits from WP2's optimal lower sets and WP6's
certificate engine.\\
WP4 holonomy & Yes, source audit & May supply a global algebraization mechanism to WP1 or
WP3.\\
WP5 conjugates & Partly & Needs normalized cocycles from the baseline; capacity calculation
can use WP2 measures.\\
WP6 certificates & Yes & Independent infrastructure; should not consume core research effort
beyond reusable tooling.\\
\bottomrule
\end{tabularx}
\end{center}

\subsection{Allocation recommendation}

The first research cycle should divide effort approximately as follows: one third on WP1,
one third split between WP2 and WP3, one sixth on the focused WP4 source audit, and the
remaining sixth on WP5 plus reusable certification.  This is not a schedule; it is a hedge
against correlated failure.  WP1 is the most direct theorem, WP2--WP3 are the most
problem-specific arithmetic experiments, and WP4--WP5 search for genuinely new sources of
global size.

\clearpage
\section{Unified lead inventory and triage}
\label{sec:inventory}

The table records all substantive clusters from the four supplied reports, including
calibrated negatives.  ``Now'' means the first experiment can be run using present objects;
``conditional'' means an explicitly named bridge must first be built; ``context'' means the
result is useful for orientation or for proving a no-go statement.

\begingroup
\scriptsize
\setlength{\LTleft}{0pt}
\setlength{\LTright}{0pt}
\setlength{\tabcolsep}{2.8pt}
\begin{longtable}{@{}L{0.225\textwidth}L{0.275\textwidth}C{0.08\textwidth}C{0.065\textwidth}L{0.30\textwidth}@{}}
\toprule
Cluster & Representative references & Source & Class & Triage / missing bridge\\
\midrule
\endfirsthead
\toprule
Cluster & Representative references & Source & Class & Triage / missing bridge\\
\midrule
\endhead
\bottomrule
\endfoot
Entire functions sharing integrality arguments & Pila I--II
\cite{PilaSharingI,PilaSharingII} & \Rfour,\Rthree & \badgeDirect & \textbf{Now.} Compute the
canonical order-two type and classify forced algebraic dependence.\\
\addlinespace
Concordant/minimal envelopes & Pila--Rodr\'iguez Villegas; Pila; Gramain
\cite{PilaRodriguezVillegas,PilaConcordant2002,Gramain1981} & \Rfour,\Rthree & \badgeAdapt &
\textbf{Now.} Verify concordance and obtain a sharp lower-growth benchmark.\\
\addlinespace
Products of two logarithms & Calegari--Dimitrov--Tang \cite{CDTProducts} &
\Rthree,\Rfour & \badgeTemplate & \textbf{Now, focused audit.} The theorem excludes the target
parameters; seek a uniform arithmetic deformation.\\
\addlinespace
Unbounded denominators / arithmetic holonomy & Calegari--Dimitrov--Tang
\cite{CDTUnbounded} & \Rfour,\Rthree & \badgeTemplate & \textbf{Conditional.} Construct an
arithmetic holonomic series attached to the orbit.\\
\addlinespace
Slope/algebraization criteria & Bost; Bost--Charles \cite{Bost2001,BostCharles} & \Rfour &
\badgeTemplate & \textbf{Conditional.} Need an algebraic formal leaf with orbit evaluation
maps.\\
\addlinespace
Multivariate fixed divisors & Prasad--Rajkumar--Reddy \cite{PrasadRajkumarReddy} & \Rone &
\badgeAdapt & \textbf{Now.} Compute on \(S_K\) with the baseline Newton polytopes.\\
\addlinespace
Total-degree multivariate factorials & Evrard \cite{Evrard2012} & \Rone & \badgeAdapt &
\textbf{Now.} Low-cost first test before the multidegree refinement.\\
\addlinespace
Composite \(b\)-orderings & Lagarias--Yangjit \cite{LagariasYangjit2024} & \Rone &
\badgeAdapt & \textbf{Now.} Test \(6,12,18,\ldots\) and weighted ideal families.\\
\addlinespace
Arbitrary-ideal, removed, and order-\(h\) orderings & Lagarias--Yangjit
\cite{LagariasYangjit2025} & \Rone & \badgeAdapt & \textbf{Now/conditional.} Compute finite
invariants; asymptotic gain must be proved.\\
\addlinespace
Adelic regular bases & Chabert--Peruginelli \cite{ChabertPeruginelli} & \Rone &
\badgeAdapt & \textbf{Conditional.} Verify residual-cardinality hypotheses for the semigroup
closures.\\
\addlinespace
Matroid structure of optimal sets & Grinberg--Petrov \cite{GrinbergPetrov} & \Rone &
\badgeAdapt & \textbf{Now.} Compute common bases and weighted simultaneous defect.\\
\addlinespace
Removed-ordering greedoids & Krachun--Mirgalimova \cite{KrachunMirgalimova} & \Rone &
\badgeAdapt & \textbf{Now.} Test robustness under deletion and boundary perturbation.\\
\addlinespace
Algorithms for refined orderings & Johnson \cite{Johnson2010} & \Rone & \badgeCompute &
\textbf{Now.} Reuse for exact finite experiments.\\
\addlinespace
Simultaneous orderings/equidistribution & Fr\k{a}czyk--Szumowicz; Byszewski et al.
\cite{FraczykSzumowicz,ByszewskiFraczykSzumowicz,SimultaneousOrderings2022} & \Rone &
\badgeTemplate & \textbf{Conditional.} Derive an energy limit for the semigroup truncations.\\
\addlinespace
Transfinite diameter and resultants & DeMarco--Rumely \cite{DeMarcoRumely} & \Rone &
\badgeTemplate & \textbf{Conditional.} Express the determinant dynamics as a polynomial-map
capacity problem.\\
\addlinespace
Adelic Fekete--Szeg\H{o} theory & Rumely \cite{RumelyFeketeSzego} & \Rone &
\badgeDiagnostic & \textbf{Context.} Determine whether local constraints are subcritical,
critical, or supercritical in capacity.\\
\addlinespace
Real analytic determinant method & Bombieri--Pila; Pila postulation
\cite{BombieriPila,PilaPostulation} & \Rfour & \badgeTemplate & \textbf{Now as design.}
Optimize staircase monomials; no direct curve theorem applies.\\
\addlinespace
\(p\)-adic/global determinant method & Heath-Brown; Salberger; Walsh
\cite{HeathBrown2002,Salberger2007,Walsh2015} & \Rthree,\Rfour & \badgeTemplate &
\textbf{Now as prototype.} Force global rank drop before taking archimedean size.\\
\addlinespace
Uniform global-field determinant bounds & Binyamini--Cluckers--Kato; CCDN
\cite{BinyaminiCluckersKato,CCDN2020} & \Rthree & \badgeTemplate & \textbf{Conditional.}
Useful after algebraization of finite orbit blocks.\\
\addlinespace
Alternants for four exponentials & Habsieger \cite{Habsieger2007} & \Rthree &
\badgeDiagnostic & \textbf{Now as audit.} Translate the current determinant and isolate the
extra hypothesis still missing.\\
\addlinespace
Modular interpolation determinants & Nesterenko \cite{Nesterenko1996} & \Rfour &
\badgeTemplate & \textbf{Context.} Demonstrates the kind of extra differential/modular
relation that can create decisive divisibility.\\
\addlinespace
House bounds & Dimitrov; Pritsker \cite{DimitrovSZ,PritskerHouse} & \Rfour & \badgeAdapt &
\textbf{Conditional.} Need low effective degree and uniform conjugate control.\\
\addlinespace
Parametric geometry of numbers & Schmidt--Summerer; Roy; DFSU
\cite{SchmidtSummerer,RoyPGN,DFSU2024} & \Rfour & \badgeTemplate & \textbf{Conditional.}
Translate integrality into a forbidden minima trajectory.\\
\addlinespace
Hypergeometric PGN / shifted logs & Kawashima--Po\"els; Hirata-Kohno et al.
\cite{KawashimaPoels,HirataKohnoMuroiWashio} & \Rfour & \badgeTemplate &
\textbf{Conditional.} Build a Pad\'e table containing \(M,A\) arithmetically.\\
\addlinespace
\(S\)-unit Pad\'e approximations & Hirata-Kohno et al. \cite{SUnitPade} & \Rfour &
\badgeTemplate & \textbf{Context/conditional.} Finiteness is insufficient; the binomial
Pad\'e construction may transfer.\\
\addlinespace
Hankel--Pad\'e irrationality criteria & Bugeaud--Han--Wen--Yao
\cite{BugeaudHanWenYao} & \Rthree & \badgeAdapt & \textbf{Now as audit.} Pair nonvanishing
with exact denominator and energy asymptotics.\\
\addlinespace
Equilibrium rational approximation & Gonchar--Rakhmanov; Stahl--Totik; Beukers
\cite{GoncharRakhmanov,StahlTotik,BeukersPade} & \Rfour & \badgeTemplate & \textbf{Now as
benchmark.} Determine whether the current residual is potential-theoretically optimal.\\
\addlinespace
Arithmetic unlikely intersections & Campagna--Dill--Wilms \cite{CampagnaDillWilms} &
\Rfour,\Rone & \badgeAdapt & \textbf{Conditional.} First force an algebraic auxiliary curve;
dimension two is also only partial.\\
\addlinespace
Effective Pfaffian counting & Binyamini--Jones--Schmidt--Thomas
\cite{BinyaminiJonesSchmidtThomas} & \Rfour & \badgeTemplate & \textbf{Context.} Effectivity
does not see denominator support.\\
\addlinespace
Analytic gcd bounds & Pasten--Wang \cite{PastenWang} & \Rfour & \badgeAdapt &
\textbf{Conditional.} Construct two algebraically independent meromorphic functions with a
forced excessive gcd counting function.\\
\addlinespace
Integer-valued definable functions & Jones--Thomas--Wilkie; Bhardwaj et al.
\cite{JonesThomasWilkie,BhardwajEtAl} & \Rfour & \badgeAdapt & \textbf{Conditional.} Need a
definable integer-valued representative on all \(\N\), not just sparse candidates.\\
\addlinespace
Algebraic unlikely intersections & BMZ; Maurin \cite{BMZ1999,Maurin2008} &
\Rtwo,\Rthree & \badgeDiagnostic & \textbf{Retire as direct route.} The orbit arc is not
algebraic.\\
\addlinespace
Complex powers & Gallinaro; Gallinaro--Kirby; Aslanyan--Kirby--Mantova
\cite{Gallinaro2025,GallinaroKirby,AslanyanKirbyMantova} & \Rfour & \badgeTemplate &
\textbf{Context.} Existence/tameness results do not give arithmetic rigidity.\\
\addlinespace
Drinfeld-module logarithms & Papanikolas; Chang--Papanikolas; Gezmi\c{s}--Namoijam
\cite{Papanikolas2008,ChangPapanikolas,GezmisNamoijam} & \Rthree & \badgeTemplate &
\textbf{Context.} Search for a characteristic-zero substitute for Frobenius difference
structure.\\
\addlinespace
Mahler relations and automatic values & Adamczewski--Faverjon
\cite{AdamczewskiFaverjon2017,AdamczewskiFaverjon2018,AdamczewskiMahlerProblem} & \Rthree &
\badgeAdapt & \textbf{Conditional.} First derive an arithmetic Mahler equation or simultaneous
automaticity.\\
\addlinespace
Digit and carry dynamics & Lagarias; Saye; Dimitrov--Howe
\cite{Lagarias2009,Saye2022,DimitrovHowe} & \Rthree & \badgeCompute & \textbf{Now for
certificates.} Infinite use needs a synchronized finite-state invariant.\\
\addlinespace
Rank matrices of logarithms & Dasgupta; Dasgupta--Kakde
\cite{DasguptaPartI,DasguptaKakdePartII} & \Rtwo,\Rthree,\Rfour & \badgeDiagnostic &
\textbf{Context.} Precisely identifies the balanced \(2\times2\) frontier.\\
\addlinespace
Analytic subgroup / period theorems & W\"ustholz; Gaudron--R\'emond
\cite{Wustholz1989,GaudronRemond} & \Rtwo,\Rfour & \badgeDiagnostic & \textbf{Conditional.}
Need an algebraic subgroup arising from integrality, not from a transcendental direction.\\
\addlinespace
Six-exponentials/algebraic-independence ceilings & Roy; Diaz
\cite{RoySixExp,Diaz1989} & \Rtwo,\Rthree & \badgeDiagnostic & \textbf{Context.} Becomes active
only after a third direction or algebraic exponent is forced.\\
\addlinespace
Explicit \(\log2,\log3\) approximation & Simons--de Weger; Simons
\cite{SimonsDeWeger,Simons2005} & \Rthree & \badgeCompute & \textbf{Now.} Improve uniform
finite certificates, not the bilinear wall.\\
\addlinespace
Hypergeometric logarithm measures & Zudilin; Marcovecchio
\cite{ZudilinSurvey,Marcovecchio2009,Marcovecchio2025} & \Rthree & \badgeCompute &
\textbf{Now for Priority H.} Verify every joint-pair constant from primary sources.\\
\addlinespace
Formal certificate interoperability & DSLean \cite{DSLean2026} & \Rthree & \badgeCompute &
\textbf{Now.} Translate Gappa/interval certificates to Lean proof terms.\\
\addlinespace
Colossally abundant application & Caveney--Nicolas--Sondow
\cite{CaveneyNicolasSondow} & \Rthree & \badgeCompute & \textbf{Downstream.} Formalize after
core definitions stabilize.\\
\addlinespace
Semiabelian gap principle & Hultberg \cite{Hultberg2026} & \Rone & \badgeTemplate &
\textbf{Context.} Requires an algebraic subvariety absent from the present arc.\\
\addlinespace
Steering-word complexity & Stephan \cite{Stephan2026} & \Rone & \badgeDiagnostic &
\textbf{Context.} Needs an exact low-complexity word forced by the counterexample.\\
\addlinespace
Transcendence measures for \(e^\alpha\) & Fischler--Rivoal
\cite{FischlerRivoal2026} & \Rone & \badgeDiagnostic & \textbf{Retire for the wall.}
\(\alpha\) is algebraic in the theorem; the target coefficients are logarithmic.\\
\addlinespace
Massold's exponential-array estimates & Massold \cite{Massold2025} & \Rone &
\badgeDiagnostic & \textbf{Already covered by baseline.} Current parameter range does not
settle the balanced \(2\times2\) case.\\
\end{longtable}
\endgroup

\subsection{Pursue, hold, and retire}

\begin{center}
\small
\begin{tabularx}{0.96\textwidth}{L{0.16\textwidth}Y}
\toprule
Bucket & Contents\\
\midrule
\textbf{Pursue now} & Pila growth audit; multivariate factorials; composite/ideal orderings;
matroid defect; global \(p\)-adic determinant prototype; CDT source audit; exact potential-
theory benchmark; finite certificate infrastructure.\\
\addlinespace
\textbf{Hold for a bridge} & House/capacity methods until conjugates are controlled;
unlikely intersections until an algebraic curve is forced; Mahler until a functional
equation is derived; o-minimal integer-valued theory until a definable representative is
constructed; analytic gcd until appropriate functions are built.\\
\addlinespace
\textbf{Use as context/no-go} & Generic BMZ/Zilber--Pink; direct analytic subgroup theorem;
ordinary Baker constant improvements; symbolic chaos without automaticity; complex-power
existence as a source of arithmetic rigidity; general capacity calculations that remain
critical.\\
\bottomrule
\end{tabularx}
\end{center}

\section{Conclusions}

The four surveys do not uncover a hidden proof, but their critical merge produces a sharper
research landscape than any report separately.

First, Pila's integer-valued-entire-function program is not merely analogous to the present
interpolation work: it is formulated on the same rank-two additive semigroup and supplies an
explicit order-two threshold.  This should be measured immediately.

Second, the generalized-factorial cluster materially extends the ordinary Bhargava analysis
already performed.  Multivariate fixed divisors, composite ideals, adelic regular bases, and
matroidal simultaneous defects are distinct invariants.  Their merit is that they can be
falsified on finite grids at the correct asymptotic scale.

Third, the CDT product-of-logarithms theorem demonstrates that the terminal type of quantity
is not universally inaccessible.  What makes its family tractable is arithmetic holonomy,
not a better generic Baker constant.  The right question is whether the candidate semigroup
supports any comparable arithmetic functional equation.

Fourth, global determinant and capacity theories clarify what the present methods are
missing.  A proof must turn local congruence into a global relation before separate
optimization, or control enough conjugates that one exponentially small embedding has global
weight.  If the corresponding capacity is exactly critical, that fact should be proved and
the route closed.

Finally, the negative results are actionable.  They rule out spending effort on generic
rational-point counting, algebraic unlikely intersections on the transcendental arc,
model-theoretic existence theorems, or linear-logarithm refinements that never engage
\cref{eq:terminal-wall}.  The surviving portfolio is smaller, more arithmetic, and equipped
with explicit decision gates.

\clearpage
\appendix
\section{Bibliography notes}

The bibliography below favors primary papers and stable arXiv records.  Several entries are
included as methodological templates rather than applicable theorems; their status is stated
in the body.  Bibliographic corrections described in \cref{sec:audit} have been incorporated.


\begingroup
\small
\sloppy
\begin{thebibliography}{99}

\bibitem{AdamczewskiFaverjon2017}
B.~Adamczewski and C.~Faverjon.
\newblock M\'ethode de Mahler: relations lin\'eaires, transcendance et applications aux nombres automatiques.
\newblock \emph{Proc. London Math. Soc. (3)} \textbf{115} (2017), 55--90.

\bibitem{AdamczewskiFaverjon2018}
B.~Adamczewski and C.~Faverjon.
\newblock M\'ethode de Mahler, transcendance et relations lin\'eaires: aspects effectifs.
\newblock \emph{J. Th\'eor. Nombres Bordeaux} \textbf{30} (2018), 557--573.

\bibitem{AdamczewskiMahlerProblem}
B.~Adamczewski.
\newblock A problem around Mahler functions.
\newblock arXiv:1303.2019.

\bibitem{AslanyanKirbyMantova}
V.~Aslanyan, J.~Kirby, and V.~Mantova.
\newblock A geometric approach to some systems of exponential equations.
\newblock \emph{Int. Math. Res. Not. IMRN} (2023), 4046--4081; arXiv:2105.12679.

\bibitem{BMZ1999}
E.~Bombieri, D.~Masser, and U.~Zannier.
\newblock Intersecting a curve with algebraic subgroups of multiplicative groups.
\newblock \emph{Int. Math. Res. Not.} (1999), 1119--1140.

\bibitem{BeukersPade}
F.~Beukers.
\newblock Pad\'e approximations in number theory.
\newblock In \emph{Pad\'e Approximation and its Applications}, Lecture Notes in Math. 888, Springer, 1981, pp.~90--99.

\bibitem{BhardwajEtAl}
N.~Bhardwaj, R.~McCulloch, N.~Ramachandran, and K.~Woo.
\newblock Integer-valued o-minimal functions.
\newblock arXiv:2404.10737.

\bibitem{BinyaminiCluckersKato}
G.~Binyamini, R.~Cluckers, and F.~Kato.
\newblock Sharp bounds for the number of rational points on algebraic curves and dimension growth, over all global fields.
\newblock \emph{Proc. London Math. Soc. (3)} \textbf{130} (2025), no.~1, e70016; arXiv:2401.03982.

\bibitem{BinyaminiJonesSchmidtThomas}
G.~Binyamini, G.~O.~Jones, H.~Schmidt, and M.~E.~M.~Thomas.
\newblock An effective Pila--Wilkie theorem for sets definable using Pfaffian functions, with some Diophantine applications.
\newblock arXiv:2301.09883.

\bibitem{BombieriPila}
E.~Bombieri and J.~Pila.
\newblock The number of integral points on arcs and ovals.
\newblock \emph{Duke Math. J.} \textbf{59} (1989), 337--357.

\bibitem{Bost2001}
J.-B.~Bost.
\newblock Algebraic leaves of algebraic foliations over number fields.
\newblock \emph{Publ. Math. Inst. Hautes \^Etudes Sci.} \textbf{93} (2001), 161--221.

\bibitem{BostCharles}
J.-B.~Bost and F.~Charles.
\newblock Quasi-projective and formal-analytic arithmetic surfaces.
\newblock arXiv:2206.14242.

\bibitem{BugeaudHanWenYao}
Y.~Bugeaud, G.-N.~Han, Z.-Y.~Wen, and J.-Y.~Yao.
\newblock Hankel determinants, Pad\'e approximations, and irrationality exponents.
\newblock \emph{Int. Math. Res. Not. IMRN} (2016), no.~5, 1467--1496; arXiv:1503.02797.

\bibitem{ByszewskiFraczykSzumowicz}
J.~Byszewski, M.~Fr\k{a}czyk, and A.~Szumowicz.
\newblock Simultaneous $p$-orderings and minimising volumes in number fields.
\newblock arXiv:1506.02696.

\bibitem{CCDN2020}
W.~Castryck, R.~Cluckers, P.~Dittmann, and K.~H.~Nguyen.
\newblock The dimension growth conjecture, polynomial in the degree and without logarithmic factors.
\newblock \emph{Algebra Number Theory} \textbf{14} (2020), no.~8, 2261--2294; arXiv:1904.13109.

\bibitem{CDTProducts}
F.~Calegari, V.~Dimitrov, and Y.~Tang.
\newblock The linear independence of $1$, $\zeta(2)$, and $L(2,\chi_{-3})$.
\newblock arXiv:2408.15403.

\bibitem{CDTUnbounded}
F.~Calegari, V.~Dimitrov, and Y.~Tang.
\newblock The unbounded denominators conjecture.
\newblock \emph{J. Amer. Math. Soc.} \textbf{38} (2025), 627--702; arXiv:2109.09040.

\bibitem{CampagnaDillWilms}
F.~Campagna, G.~A.~Dill, and R.~Wilms.
\newblock Arithmetic unlikely intersections in powers of the multiplicative group.
\newblock arXiv:2607.15741.

\bibitem{CaveneyNicolasSondow}
G.~Caveney, J.-L.~Nicolas, and J.~Sondow.
\newblock On SA, CA, and GA numbers.
\newblock \emph{Ramanujan J.} \textbf{29} (2012), 359--384.

\bibitem{ChabertPeruginelli}
J.-L.~Chabert and G.~Peruginelli.
\newblock Adelic versions of the Weierstrass approximation theorem.
\newblock \emph{J. Pure Appl. Algebra} \textbf{222} (2018); arXiv:1511.03465.

\bibitem{ChangPapanikolas}
C.-Y.~Chang and M.~A.~Papanikolas.
\newblock Algebraic independence of periods and logarithms of Drinfeld modules.
\newblock \emph{J. Amer. Math. Soc.} \textbf{25} (2012), 123--150.

\bibitem{DFSU2024}
T.~Das, L.~Fishman, D.~Simmons, and M.~Urba\'nski.
\newblock A variational principle in the parametric geometry of numbers.
\newblock \emph{Adv. Math.} \textbf{437} (2024), 109435.

\bibitem{DSLean2026}
T.~Rowney, R.~Ahuja, J.~Avigad, and S.~Welleck.
\newblock DSLean: a framework for type-correct interoperability between Lean~4 and external DSLs.
\newblock arXiv:2602.18657.

\bibitem{DasguptaKakdePartII}
S.~Dasgupta and M.~Kakde.
\newblock Ranks of matrices of logarithms of algebraic numbers II: the Matrix Coefficient Conjecture.
\newblock arXiv:2408.08178.

\bibitem{DasguptaPartI}
S.~Dasgupta.
\newblock Ranks of matrices of logarithms of algebraic numbers I: the theorems of Baker and Waldschmidt--Masser.
\newblock arXiv:2303.02037.

\bibitem{DeMarcoRumely}
L.~DeMarco and R.~Rumely.
\newblock Transfinite diameter and the resultant.
\newblock arXiv:math/0601109; published 2007.

\bibitem{Diaz1989}
G.~Diaz.
\newblock Grands degr\'es de transcendance pour des familles d'exponentielles.
\newblock \emph{J. Number Theory} \textbf{31} (1989), no.~1, 1--23.

\bibitem{DimitrovHowe}
V.~Dimitrov and E.~W.~Howe.
\newblock Powers of $3$ with few nonzero bits and a conjecture of Erd\H{o}s.
\newblock arXiv:2105.06440.

\bibitem{DimitrovSZ}
V.~Dimitrov.
\newblock A proof of the Schinzel--Zassenhaus conjecture on polynomials.
\newblock arXiv:1912.12545.

\bibitem{Evrard2012}
S.~Evrard.
\newblock Bhargava's factorials in several variables.
\newblock \emph{J. Algebra} \textbf{372} (2012), 134--148.

\bibitem{FischlerRivoal2026}
S.~Fischler and T.~Rivoal.
\newblock A new transcendence measure for the values of the exponential function at algebraic arguments.
\newblock arXiv:2502.17992.

\bibitem{FraczykSzumowicz}
M.~Fr\k{a}czyk and A.~Szumowicz.
\newblock On the optimal rate of equidistribution in number fields.
\newblock arXiv:1810.11110.

\bibitem{Gallinaro2025}
F.~Gallinaro.
\newblock Solving systems of equations of raising-to-powers type.
\newblock \emph{Israel J. Math.} \textbf{270} (2025), 229--256; arXiv:2103.15675.

\bibitem{GallinaroKirby}
F.~Gallinaro and J.~Kirby.
\newblock Quasiminimality of complex powers.
\newblock arXiv:2304.06450.

\bibitem{GaudronRemond}
\'E.~Gaudron and G.~R\'emond.
\newblock Th\'eor\`eme des p\'eriodes et degr\'es minimaux d'isog\'enies.
\newblock \emph{Comment. Math. Helv.} \textbf{89} (2014), 343--403; arXiv:1105.1230.

\bibitem{GezmisNamoijam}
O.~Gezmi\c{s} and C.~Namoijam.
\newblock On the algebraic independence of logarithms of Anderson $t$-modules.
\newblock arXiv:2407.18916.

\bibitem{GoncharRakhmanov}
A.~A.~Gonchar and E.~A.~Rakhmanov.
\newblock Equilibrium distributions and degree of rational approximation of analytic functions.
\newblock \emph{Math. USSR Sb.} \textbf{62} (1989), 305--348.

\bibitem{Gramain1981}
F.~Gramain.
\newblock Sur le th\'eor\`eme de Fukasawa--Gel'fond.
\newblock \emph{Invent. Math.} \textbf{63} (1981), 495--506.

\bibitem{GrinbergPetrov}
D.~Grinberg and F.~Petrov.
\newblock A greedoid and a matroid inspired by Bhargava's $p$-orderings.
\newblock \emph{Electron. J. Combin.} \textbf{28} (2021), no.~3, Paper 3.6; arXiv:1909.01965.

\bibitem{Habsieger2007}
L.~Habsieger.
\newblock A few remarks related to the four exponentials conjecture.
\newblock \emph{Publ. Math. Debrecen} \textbf{70} (2007), no.~1--2, 179--194.

\bibitem{HeathBrown2002}
D.~R.~Heath-Brown.
\newblock The density of rational points on curves and surfaces.
\newblock \emph{Ann. of Math. (2)} \textbf{155} (2002), 553--595.

\bibitem{HirataKohnoMuroiWashio}
N.~Hirata-Kohno, T.~Muroi, and M.~Washio.
\newblock Approximation measures for shifted logarithms of algebraic numbers via effective Poincar\'e--Perron.
\newblock arXiv:2406.10820.

\bibitem{Hultberg2026}
N.~Hultberg.
\newblock New gap principle for semiabelian varieties using globally valued fields.
\newblock arXiv:2601.04972.

\bibitem{Johnson2010}
K.~Johnson.
\newblock Computing $r$-removed $P$-orderings and $P$-orderings of order $h$.
\newblock \emph{Actes des rencontres du CIRM} \textbf{2} (2010), no.~2, 33--40.

\bibitem{JonesThomasWilkie}
G.~O.~Jones, M.~E.~M.~Thomas, and A.~J.~Wilkie.
\newblock Integer-valued definable functions.
\newblock \emph{Bull. London Math. Soc.} \textbf{44} (2012), 1285--1291.

\bibitem{KawashimaPoels}
M.~Kawashima and A.~Po\"els.
\newblock Pad\'e approximation for a class of hypergeometric functions and parametric geometry of numbers.
\newblock arXiv:2202.10782.

\bibitem{KrachunMirgalimova}
D.~Krachun and A.~Mirgalimova.
\newblock Strong greedoid structure of $r$-removed $P$-orderings.
\newblock \emph{Electron. J. Combin.} \textbf{31} (2024), no.~4; arXiv:2309.10104.

\bibitem{Lagarias2009}
J.~C.~Lagarias.
\newblock Ternary expansions of powers of $2$.
\newblock \emph{J. London Math. Soc. (2)} \textbf{79} (2009), 562--588; arXiv:math/0512006.

\bibitem{LagariasYangjit2024}
J.~C.~Lagarias and W.~Yangjit.
\newblock The factorial function and generalizations, extended.
\newblock \emph{J. Number Theory} \textbf{259} (2024), 131--170; arXiv:2310.12949.

\bibitem{LagariasYangjit2025}
J.~C.~Lagarias and W.~Yangjit.
\newblock $B$-orderings for all ideals $B$ of Dedekind domains and generalized factorials.
\newblock arXiv:2502.19072.

\bibitem{Marcovecchio2009}
R.~Marcovecchio.
\newblock The Rhin--Viola method for $\log 2$.
\newblock \emph{Acta Arith.} \textbf{139} (2009), no.~2, 147--184.

\bibitem{Marcovecchio2025}
R.~Marcovecchio.
\newblock Multiple Legendre polynomials in Diophantine approximation.
\newblock arXiv:2512.12890.

\bibitem{Massold2025}
H.~Massold.
\newblock Transcendence degrees of fields generated by exponentials of products.
\newblock arXiv:2506.01123.

\bibitem{Maurin2008}
G.~Maurin.
\newblock Courbes alg\'ebriques et \`equations multiplicatives.
\newblock \emph{Math. Ann.} \textbf{341} (2008), 789--824.

\bibitem{Nesterenko1996}
Yu.~V.~Nesterenko.
\newblock Modular functions and transcendence questions.
\newblock \emph{Sb. Math.} \textbf{187} (1996), 1319--1348.

\bibitem{Papanikolas2008}
M.~A.~Papanikolas.
\newblock Tannakian duality for Anderson--Drinfeld motives and algebraic independence of Carlitz logarithms.
\newblock \emph{Invent. Math.} \textbf{171} (2008), 123--174.

\bibitem{PastenWang}
H.~Pasten and J.~Wang.
\newblock GCD bounds for analytic functions.
\newblock \emph{Int. Math. Res. Not. IMRN} (2017), no.~1, 47--95.

\bibitem{PilaConcordant2002}
J.~Pila.
\newblock Concordant sequences and concordant entire functions.
\newblock \emph{Rend. Circ. Mat. Palermo (2)} \textbf{51} (2002), 51--82.

\bibitem{PilaPostulation}
J.~Pila.
\newblock Geometric and arithmetic postulation of the exponential function.
\newblock \emph{J. Austral. Math. Soc. Ser. A} \textbf{54} (1993), no.~1, 111--127.

\bibitem{PilaRodriguezVillegas}
J.~Pila and F.~Rodr\'iguez Villegas.
\newblock Concordant sequences and integral-valued entire functions.
\newblock \emph{Acta Arith.} \textbf{88} (1999), 239--268.

\bibitem{PilaSharingI}
J.~Pila.
\newblock Entire functions sharing arguments of integrality, I.
\newblock \emph{Int. J. Number Theory} \textbf{5} (2009), no.~2, 339--353.

\bibitem{PilaSharingII}
J.~Pila.
\newblock Entire functions sharing arguments of integrality, II.
\newblock \emph{Publ. Math. Debrecen} \textbf{73} (2008), no.~1--2, 215--231.

\bibitem{PrasadRajkumarReddy}
D.~Prasad Semwal, K.~Rajkumar, and A.~S.~Reddy.
\newblock Fixed divisor of a multivariate polynomial and generalized factorials in several variables.
\newblock \emph{J. Korean Math. Soc.} \textbf{55} (2018), 1305--1320; arXiv:1803.05780.

\bibitem{PritskerHouse}
I.~E.~Pritsker.
\newblock House of algebraic integers symmetric about the unit circle.
\newblock arXiv:2101.06710.

\bibitem{RoyPGN}
D.~Roy.
\newblock On Schmidt and Summerer parametric geometry of numbers.
\newblock \emph{Ann. of Math. (2)} \textbf{182} (2015), 739--786.

\bibitem{RoySixExp}
D.~Roy.
\newblock Matrices whose coefficients are linear forms in logarithms.
\newblock \emph{J. Number Theory} \textbf{41} (1992), no.~1, 22--47.

\bibitem{RumelyFeketeSzego}
R.~Rumely.
\newblock The Fekete--Szeg\H{o} theorem with local rationality conditions on curves.
\newblock arXiv:1203.1213.

\bibitem{SUnitPade}
N.~Hirata-Kohno, M.~Kawashima, A.~Po\"els, and M.~Washio.
\newblock $S$-unit equation in two variables and Pad\'e approximations.
\newblock arXiv:2211.14399.

\bibitem{Salberger2007}
P.~Salberger.
\newblock On the density of rational and integral points on algebraic varieties.
\newblock \emph{J. Reine Angew. Math.} \textbf{606} (2007), 123--147.

\bibitem{Saye2022}
R.~I.~Saye.
\newblock On two conjectures concerning the ternary digits of powers of two.
\newblock \emph{J. Integer Seq.} \textbf{25} (2022); arXiv:2202.13256.

\bibitem{SchmidtSummerer}
W.~M.~Schmidt and L.~Summerer.
\newblock Parametric geometry of numbers and applications.
\newblock \emph{Acta Arith.} \textbf{140} (2009), 67--91.

\bibitem{Simons2005}
J.~Simons.
\newblock A theorem on the nonexistence of $2$-cycles for the $3x+1$ problem.
\newblock \emph{Math. Comp.} \textbf{74} (2005), 1565--1572.

\bibitem{SimonsDeWeger}
J.~Simons and B.~de~Weger.
\newblock Theoretical and computational bounds for $m$-cycles of the $3n+1$ problem.
\newblock \emph{Acta Arith.} \textbf{117} (2005), no.~1, 51--70.

\bibitem{SimultaneousOrderings2022}
A.~Szumowicz.
\newblock Simultaneous $p$-orderings and equidistribution.
\newblock arXiv:2207.08233.

\bibitem{StahlTotik}
H.~Stahl and V.~Totik.
\newblock \emph{General Orthogonal Polynomials}.
\newblock Encyclopedia of Mathematics and its Applications 43, Cambridge University Press, 1992.

\bibitem{Stephan2026}
R.~Stephan.
\newblock Superlinear complexity of the $(3/2)^n$ steering word.
\newblock arXiv:2607.11648.

\bibitem{WaldschmidtIntegerValued}
M.~Waldschmidt.
\newblock Integer-valued functions, Hurwitz functions and related topics.
\newblock arXiv:2002.01223.

\bibitem{Walsh2015}
M.~N.~Walsh.
\newblock Bounded rational points on curves.
\newblock \emph{Int. Math. Res. Not. IMRN} (2015), 5644--5658.

\bibitem{Wustholz1989}
G.~W\"ustholz.
\newblock Algebraische Punkte auf analytischen Untergruppen algebraischer Gruppen.
\newblock \emph{Ann. of Math. (2)} \textbf{129} (1989), 501--517.

\bibitem{ZudilinSurvey}
W.~Zudilin.
\newblock An essay on irrationality measures of $\pi$ and other logarithms.
\newblock arXiv:math/0404523.

\end{thebibliography}
\endgroup

\end{document}
